\documentclass[11pt]{article}

% ---------------------------------------------------------------------
% Preamble: standard TeX Live packages only.
% ---------------------------------------------------------------------
\usepackage[T1]{fontenc}
\usepackage[utf8]{inputenc}
\usepackage{lmodern}
\usepackage[margin=1in]{geometry}
\usepackage{graphicx}
\usepackage{booktabs}
\usepackage{array}
\usepackage{amsmath}
\usepackage{amssymb}
\usepackage{xcolor}
\usepackage{url}
\usepackage[colorlinks=true, linkcolor=blue!50!black, citecolor=blue!50!black, urlcolor=blue!50!black, pdftitle={Two Routes to 2D Animation on an 8 GB Consumer GPU}, pdfauthor={Author Name (TODO: fill in before submission)}, pdfsubject={2D animation generation on consumer hardware}, pdfkeywords={2D animation, consumer GPU, evaluation methodology, small language models, ControlNet, AnimateDiff, decoding artefacts}]{hyperref}

% Wrapped-text column types for wide tables.
\newcolumntype{L}[1]{>{\raggedright\arraybackslash}p{#1}}
\newcolumntype{C}[1]{>{\centering\arraybackslash}p{#1}}

% Section symbol shorthand used throughout for in-text cross-references;
% since every \section/\subsection below is numbered sequentially exactly
% as in the Markdown source (paper/PAPER.md), a literal "\S4.3" reproduces
% the source's own "\S4.3" cross-references without a separate \label/\ref
% network.
\newcommand{\secsym}{\S}

\title{Two Routes to 2D Animation on an 8~GB Consumer GPU:\\
A Closed Route, an Open One, and What a Small Model\\
Can and Cannot Learn to Write}

% TODO: no author name, affiliation, or contact email is recorded anywhere
% in the project's own documents (ATTEMPTS.md, PAPER.md, RESEARCH.md,
% RESEARCH2.md). Fill this in before compiling for submission -- do not
% submit with a placeholder author.
\author{Author Name \\ \small (affiliation to be added) \\ \small \texttt{email@example.com}}
\date{}

\begin{document}

\maketitle

\begin{quote}
\small\itshape
\textbf{Status:} working draft, current through Attempt 23
(\texttt{ATTEMPTS.md}, 2929 lines). All numbers in this draft are taken
from \texttt{ATTEMPTS.md}, \texttt{critic/README.md}, and
\texttt{paper/RESEARCH2.md} as they stand at time of writing, with an
inline source note on every table; \texttt{paper/RESEARCH.md} and
\texttt{paper/RESEARCH2.md} are cited for related-work claims only. Every
superseded number from earlier evaluation runs has been replaced with the
corrected figure --- see \secsym5 for what changed and why.
\end{quote}

\begin{abstract}
We report twenty-three logged attempts (\texttt{ATTEMPTS.md}) to build a
small, locally-trainable system for 2D animation generation on a single
8~GB consumer laptop GPU (RTX~4060 Laptop, thermally constrained to a
documented 50\% duty cycle), together with the evaluation harness built to
judge the results honestly. The project's arc has two halves. The first
(Attempts 17--20) pursues styled character animation through Stable
Diffusion~1.5 + ControlNet + AnimateDiff, and closes that route by
measurement: even after correcting a genuine skeleton-topology bug and
isolating the deeper cause (a stylized character's proportions sit outside
a human-pose model's training distribution), and even after adding the
first mechanism in the project's history to reduce flicker at all
($36.2\times \rightarrow 25.0\times$, a real 31\% cut), every configuration
tried remains a hard failure ($25.0\times$ against an $8\times$ bar) and
the fix that helps flicker costs pose obedience outright (an obeyed pose
in 1 of 4 frames drops to 0 of 8). The second half (Attempts 21--23)
changes direction entirely: a CPU-only, symbolic scene-script compositor
produces the project's first multi-character animation --- three and four
independently-acted characters with background, props, occlusion and a
persistent ball --- and clears the evaluation harness's temporal gates by
roughly $25\times$ the margin the diffusion route's best configuration
needed (warp error 0.0055 against a 0.06 pass bar, versus 0.135 for
AnimateDiff). A 60M-parameter model then learns to \emph{write} these
scripts from natural language, training in under nine minutes on CPU with
no GPU at any point: it parses and renders 100\% of held-out prompts,
satisfies every stated fact in 90.7\% of them in-distribution, and
generalizes to unseen phrasing perfectly at the format level --- but
suffers a 65.7-point compositional generalization gap on cast size
specifically. A controlled probe traces this exactly: emitted cast size
equals clause count, not the stated numeral, in 15 of 15 trials, and the
cause is a confound the project introduced into its own training data
(clause count equalled cast size in 99.9\% of synthesized examples).
Deconfounding the data and retraining --- same model, same
hyperparameters --- cuts the gap to 23.5 points at a measured 4-point
in-distribution cost.

A third attempt on this half of the arc produces the paper's headline
finding: \textbf{the ``poorly staged'' scenes the project attributed to
the model are mostly a decoding artefact, not a model defect.} Every
position number reported for Attempts 21--22 used beam search, which
returns the approximately most-likely \emph{sequence} --- for a slot the
prompt never constrains, that is by definition the marginal mode. Decoding
the \emph{same checkpoint} on the \emph{same prompts} four ways and
scoring each with a new staging metric (independently calibrated: the
project's two hand-written scenes score 0.912 mean with zero overlap
against eleven model-written scenes at 0.352 mean) shows staging rising
from 0.188 to 0.862 under sampling instead of beam search on one
checkpoint, and 0.333 to 0.845 on another --- same weights, same prompts,
only the decoding argument changed. Thirty scenes under beam search emit
three distinct $x$-coordinates in total; the same weights, sampled, emit
fifty-five. A companion ablation shows what would have shipped without
this check: under beam search, a deterministic post-decode spacing pass
does essentially all the visible work (staging $0.178 \rightarrow 0.587$),
which reads as a successful engineering fix and is, by the project's own
later measurement, the wrong scientific claim --- sampling alone, no pass,
reaches 0.879. We state this exactly: \emph{good engineering, bad science
claim.} Sample-and-rerank then recovers what sampling costs in structural
validity ($60.7\% \rightarrow 94.0\%$ in-distribution, the highest this
project has recorded) while leaving staging untouched, and the resulting
pipeline produces two deliverables that clear the evaluation harness ---
one that stages above the worse of the two hand-written reference clips
for the first time in this project's history, and one with six
independently-animated characters, the project's highest measured motion
coverage. This is not reported as an unqualified win: the same attempt's
own control shows the model's apparent obedience to stated spatial
instructions (e.g.\ ``leftmost'') is within noise of a mismatched-prompt
baseline (+2.7 points in distribution, a falsified prediction stated as
such), sampling costs 13--20 points of action-fidelity and layout
obedience that reranking does not recover, and cast size now saturates at
nine rather than six --- the wall relocated by exactly the width of the
training-range extension, not removed, which we report as confirming a
known general limitation of sequence-to-sequence counting rather than a
defect specific to this system.

We report three central findings, not one or two. First, extending
Attempts 17--20: style, motion, identity and pose behave as separable
defects that trade against each other on this hardware and model family
rather than compounding for free. Second, from Attempt 22: a small model
trained on synthesized data can look incapable of a whole class of
reasoning (here, counting) when the truer diagnosis is that its training
data never asked it to do that reasoning at all --- checkable by a
controlled probe, rather than assumed. Third, from Attempt 23: an apparent
model defect can itself be a decoding artefact, checkable the same way ---
by holding the weights and the prompts fixed and varying only the decoding
strategy --- and this is the project's third corrected diagnosis of this
kind, after the COCO-18 control that overturned a topology explanation
(Attempt 19) and the clause-count probe that overturned ``the model cannot
count'' (Attempt 22); we treat the pattern of finding, and then
re-examining, its own apparent failures as a methodological contribution
in its own right. We also report, as a contribution in its own right, four
cases in which this project's own evaluation harness or decoding pipeline
produced a false or unmeasurable verdict and was corrected in the open: a
CLIP-based content-correctness metric that cannot separate known-bad
geometry from correct geometry in this domain, an identity metric that
silently mixed two measurement bases mid-clip and cost one real clip its
verdict twice, pose/person detectors that read zero on every stick figure
this project has produced, and beam search collapsing an apparently-fine
learned distribution into a single unrepresentative point. Finally we
answer the project's own standing indictment --- ``the models never beat
the script'' --- directly and without spin. Attempt 22's model-written
scene is thinner (2 characters, 2 props, 4 events) than the hand-written
one (3 characters, 7 props, 8 events) and both are USABLE; that comparison
still holds. What has changed since is narrower and precise: with a
corrected decoder, a model-written scene now stages \emph{better} than the
weaker of the two hand-written reference clips ($0.886$ vs.\ $0.825$) ---
composition, not direction, since the same attempt's own control shows the
model still does not reliably obey a stated left/right instruction.
``100\% hand-written'' is no longer true of the project's best output
category; ``the model takes direction'' is not yet true of it either, and
we say both.
\end{abstract}

\section{Introduction}
\label{sec:intro}

Most published work on video and animation generation assumes tens to
hundreds of gigabytes of accelerator memory and training budgets far
beyond a single consumer machine. This paper asks a narrower, more
answerable question: \textbf{what is actually reachable in 2D animation
generation on one 8~GB laptop GPU with no cloud budget, and what fails,
generalizes, or partially works, and why?} We treat the answer as worth
publishing in its own right --- not as a preamble to a working system, but
as the deliverable.

The project (\texttt{2DVideoGen}) logs twenty-three attempts. The first
four eras (equation-based rendering, distillation, real pose/motion data,
and a symbolic vector format) establish the representational groundwork
and are summarized in \secsym3. The paper's results then split into two
connected arcs, and its central claims are drawn from both.

\textbf{Arc one (Attempts 17--20) closes a route by measurement.} A hybrid
of a small language model for motion and Stable Diffusion~1.5 + ControlNet
+ AnimateDiff for style is built, measured, diagnosed twice (once wrongly,
once correctly), and ultimately shown to fail its own evaluation gate
under every configuration tried --- including the one configuration that
delivered this project's first genuine reduction in flicker. The specific
finding that survives this arc, and that we promote to the paper's first
central claim, is that flicker, identity and pose obedience behave as
three separable defects on this hardware and model family: fixing one
measurably costs another, rather than the three compounding toward a
shared solution as more mechanism is added.

\textbf{Arc two (Attempts 21--23) opens a different route and gets
further on it than anything earlier in the project.} Abandoning
diffusion-rendered style entirely, a CPU-only symbolic scene compositor
produces this project's first-ever multi-character animation --- several
independently acted stick figures, a background, props, correct
occlusion, and a stateful ball --- and clears the evaluation harness's
temporal-stability gates by a wide margin. A 60M-parameter language model
is then trained, in under nine minutes on a CPU, to write these scripts
from natural language. It generalizes at the format level essentially
perfectly and fails compositionally in one specific, diagnosable way: a
spurious correlation the project introduced into its own synthesized
training data. Isolating that confound with a controlled probe, fixing it
at the data level, and reporting the resulting trade honestly is the
paper's second central claim, and we hold it up as a general
methodological point about evaluating small models trained on synthesized
data --- a model can look incapable of a whole class of reasoning when the
truer fault is what its training data ever asked it to learn.

\textbf{A third attempt on arc two (Attempt 23) then reopens the diagnosis
of arc two's own most-cited defect.} Attempt 22 named ``staging'' ---
characters clustering at nearly the same position --- as the single
biggest remaining defect and attributed it to unsupervised position
slots. Attempt 23 builds an independently-calibrated staging metric,
decodes the \emph{identical checkpoint} on the \emph{identical prompts}
under four decoding strategies, and finds the defect is mostly a property
of beam search, not of what the model had learned: staging rises from
0.188 to 0.862 on one checkpoint and 0.333 to 0.845 on another purely by
sampling instead of beam-searching. A companion ablation shows a
deterministic spacing pass would otherwise have been mistaken for the
fix, when it is doing most of its visible work only because beam search
was suppressing an already-adequate learned distribution --- reported as
\emph{good engineering, bad science claim}. This is the paper's third
central finding, and it completes a pattern: this project has now
corrected its own diagnosis of an apparent model failure three times (the
COCO-18 topology control in Attempt 19, the clause-count probe in Attempt
22, and the decoder ablation in Attempt 23), each time by building a
control that could have shown the opposite result and did not assume the
first plausible explanation. We treat that pattern itself as a
methodological contribution (\secsym5).

Running through all three arcs is a fourth, cross-cutting contribution:
this project's own evaluation harness and decoding pipeline produced at
least four false or unmeasurable verdicts in the course of the work, each
caught, diagnosed, and corrected in the open rather than quietly fixed. We
treat that as a methodological contribution in its own right (\secsym5),
because the failure mode --- a metric or decoding default silently
mismeasuring or misrepresenting non-photographic, stylized, symbolic, or
unsupervised-slot content --- is not specific to this project's tooling.

\subsection*{Contributions}

\begin{enumerate}
\item \textbf{A measured demonstration that flicker, identity and pose
  fidelity are three separate defects that trade against each other on
  this hardware and model family, not three sub-problems of one shared
  deficiency.} An explicit temporal-attention module (AnimateDiff)
  delivers this project's first genuine reduction in flicker ---
  $36.2\times \rightarrow 25.0\times$, a 31\% cut in the styled term
  against an unchanged guide --- while simultaneously erasing pose
  obedience (an obeyed guide pose in 1 of 4 frames drops to 0 of 8), and it
  still leaves flicker roughly $3\times$ over the evaluation harness's own
  $8\times$ bar. An identity anchor (IP-Adapter) that fixes the separate
  identity-preservation failure (worst-case similarity $0.744 \rightarrow
  0.896$) costs a further, small increase in flicker. Neither mechanism
  helps the other's number.

\item \textbf{A multi-character symbolic animation system that runs
  entirely on CPU and clears the project's own evaluation gates by a wide
  margin.} \texttt{scenescript.py} composes a cast, background, props and
  a timeline into a rendered video with correct occlusion and a
  persistent, stateful prop; the resulting clip passes warp error at
  0.0055 against a 0.06 pass bar --- roughly $25\times$ under the
  hard-fail line where the best diffusion configuration (Attempt 20) sat
  at 0.135 --- and a script written after the compositor's code was frozen
  renders correctly on the first attempt with zero code changes.

\item \textbf{A 60M-parameter model that learns to write these scripts
  from natural language, training in under nine minutes on CPU with no GPU
  used at any point}, generalizing to unseen phrasing at the format level
  perfectly (100\% parse, 100\% render, both in- and out-of-distribution)
  while showing a specific, diagnosable compositional failure on cast
  size.

\item \textbf{A controlled probe that traces a 65.7-point generalization
  gap to a single spurious correlation the project put into its own
  training data}, and a methodological point we generalize beyond this
  project: a model that appears unable to perform a class of reasoning
  (counting, here) should be tested with a probe that holds the
  reasoning's inputs independent before concluding the model lacks the
  capability --- in this case the model could produce any cast size, it
  had simply learned to read clause count (99.9\% correlated with cast
  size in the original synthesized data) instead of the stated numeral.
  Deconfounding the data and retraining, with no other change, cut the gap
  to 23.5 points at a measured, honestly reported 4-point in-distribution
  cost.

\item \textbf{A decoding-strategy ablation that reattributes the
  project's own most-cited model defect.} Attempt 22 attributed poor
  staging (characters clustering near-identically) to the model. Decoding
  the \emph{same checkpoint} on the \emph{same prompts} under beam search
  versus sampling shows staging rising from 0.188 to 0.862 on one
  checkpoint and 0.333 to 0.845 on another --- an artefact of beam search
  reporting the mode of an already-adequate learned distribution, not a
  defect in what the model had learned. A companion four-cell ablation
  shows a deterministic post-decode spacing pass, which would otherwise
  have been reported as the fix, is doing almost all of its visible work
  only because beam search was suppressing the same distribution sampling
  reveals (staging $0.178 \rightarrow 0.587$ via the pass versus $0.178
  \rightarrow 0.879$ via decoding alone) --- stated exactly as the project
  states it: \emph{good engineering, bad science claim.} This is the
  paper's third corrected diagnosis of an apparent model failure (after
  the COCO-18 topology control, contribution 7, and the clause-count
  probe, contribution 4), and we treat the pattern of self-correction as a
  contribution distinct from any one of the three individual corrections.

\item \textbf{A direct, unspun answer to the project's own standing
  indictment} (``the models never beat the script''): Attempt 22's
  model-written scene remains measurably thinner than the hand-written one
  (2 characters, 2 props, 4 events, against 3, 7 and 8) though both clear
  the evaluation harness's USABLE bar; with the decoding fix, a
  model-written scene now \emph{stages} above the weaker of the two
  hand-written reference clips ($0.886$ vs.\ $0.825$) for the first time,
  while the same attempt's own control shows the model still does not
  reliably obey a stated left/right instruction --- composition improved
  measurably, direction did not. ``100\% hand-written'' is no longer true
  of the project's best output category; the model does not yet take
  direction, and we say both plainly.

\item \textbf{Four corrected evaluation or decoding failures, reported as
  a methodological contribution in their own right}: a CLIP-embedding
  content-correctness metric that cannot separate known-bad geometry from
  correct geometry in this domain; an identity-preservation metric that
  silently compared crops taken under two different measurement regimes
  within one clip, costing one real clip its verdict, discovered a second
  time retroactively once the first fix was applied; pose/person detectors
  (CMU OpenPose, YOLO11-pose) that read zero or near-zero on every stick
  figure this project has produced, hand-written or model-generated,
  independent of cast size; and beam search collapsing an adequate learned
  position distribution into a single unrepresentative point, discovered
  only because a control decoded the same weights a second way. We
  generalize this: both metrics built against photographic training data
  and decoding defaults tuned for fluent text need independent
  verification, not assumption, before being trusted on stylized,
  symbolic, or unsupervised-slot generation.

\item \textbf{A corrected two-stage diagnosis of a pose-conditioning
  failure in the diffusion route.} A ControlNet-OpenPose hybrid ignored an
  exact, hand-verified pose guide; the first hypothesis (wrong keypoint
  topology) was fixed completely and the pose was \emph{still} ignored. A
  controlled ablation against a genuine photo annotation and a hand-built
  canonical skeleton isolated the real cause: the stylized character's
  metric proportions, not its topology, put the guide outside the
  conditioning model's training distribution --- a finding we state as
  general to human-pose conditioning of non-human-proportioned characters,
  not specific to this implementation.

\item \textbf{The first stacked VRAM measurement of a ControlNet +
  AnimateDiff motion-module pipeline on an 8~GB consumer card}, closing an
  open question the project's own literature survey (\texttt{paper/RESEARCH.md})
  could only mark unverified: 4.41~GiB resident weights, 6.15~GiB peak
  reserved at $512\times8$, fitting without offload; only $512\times16$
  requires it.

\item \textbf{A full record across four earlier representational eras}
  (equations, distilled nets, real pose data, symbolic vector graphics),
  each capped for a different, measured reason, kept so the dead ends are
  not silently re-attempted.
\end{enumerate}

We do not claim a fully solved system. Arc one is a closed route, stated
as closed. Arc two produces the project's best output to date ---
including a scene that stages better than one of the project's two
hand-written reference clips --- and still falls short of a well-staged,
prompt-controllable, natural-language-driven animation system: the model
does not obey stated spatial instructions above the base rate, cast size
saturates at nine, and sampling's staging gain costs 13--20 points of
action fidelity that reranking does not recover. \secsym13 says plainly
what remains, including what this paper does not evaluate at all (no user
study, no comparison against a competing system).

\section{Related Work}
\label{sec:relwork}

We draw only on citations that appear in this project's own literature
surveys (\texttt{paper/RESEARCH.md}, covering Attempts 17--20's diffusion
route, and \texttt{paper/RESEARCH2.md}, written in parallel with Attempt
23 to cover the staging/counting/constrained-decoding literature for the
scene-script route); we do not introduce new references beyond what these
two surveys cite. Attempt 21 is pure systems work internal to the project
and cites no external literature in \texttt{ATTEMPTS.md}; we do not invent
citations for it.

\paragraph{Temporal consistency for per-frame diffusion.} AnimateDiff
\cite{guo2023animatediff} inserts a motion module, trained separately,
into a frozen Stable Diffusion~1.5 UNet, and transfers to any SD1.5
checkpoint without retraining; community-reported VRAM for
$512\times512\times16$-frame generation ranges roughly 8--21~GB depending
on settings, i.e.\ workable on an 8~GB card only at reduced
resolution/frame count. AnimateDiff-Lightning
\cite{lin2024animatedifflightning} distills the same motion module to
1--8-step inference for roughly $10\times$ speedup; its VRAM footprint
relative to base AnimateDiff is not separately quantified in the sources
surveyed. Training-free alternatives include Text2Video-Zero
\cite{khachatryan2023text2videozero}, which anchors appearance via
cross-frame attention but is reported at a 12~GB minimum --- above this
project's budget as published --- and FRESCO \cite{yang2024fresco}, a
zero-shot spatial- and temporal-correspondence method superseding
Rerender-A-Video \cite{yang2023rerender} by the same authors' own account.
TokenFlow \cite{geyer2023tokenflow} propagates diffusion features via
inter-frame correspondence rather than warping pixels, but is reported at
5091~s for 240 frames on a 24~GB card (PipeFlow \cite{pipeflow2025}, which
is $9.6\times$ faster at the same task) --- a wall-clock cost that rules
it out for iterative work on a laptop already duty-cycled for thermal
reasons. ControlVideo's full cross-frame attention is the most expensive
of this group by construction (quadratic in frame count) and was not
adopted without local measurement. Newer sparse/lightweight
temporal-attention methods (Sparse Forcing \cite{sparseforcing2026},
LiteAttention \cite{liteattention2025}, Sparse VideoGen
\cite{sparsevideogen2025}, HyperVAttention \cite{hypervattention2026})
target diffusion-transformer video backbones rather than the SD1.5+UNet
stack used here and are not shown composing with ControlNet-SD1.5 in the
sources surveyed.

\paragraph{Identity preservation.} AnimateAnyone and MagicAnimate (cited
via Animate-X \cite{animatex2024} and a MagicAnimate ResearchGate record
\cite{magicanimate2023}) add a ReferenceNet --- a duplicated UNet branch
injecting appearance features from one fixed reference image --- plus a
lightweight Pose Guider, addressing exactly the ``every frame is a
different character'' failure mode reported in this project's own Attempt
18 (\secsym7.1). MagicAnimate is reported to still show fine-grained
inconsistency and background drift under its own guidance signal.
IP-Adapter / IP-Adapter-FaceID is a lighter-weight alternative (frozen
base, trained cross-attention adapter only), reported at roughly 6--8~GB
for SD1.5 at $512\times512$ --- the only concrete figure found that sits
inside this project's budget without qualification, though it is a weaker
identity lock than a full ReferenceNet and its full-body/clothing
consistency is untested in the sources surveyed.

\paragraph{Script-to-multi-shot generation.} VideoDirectorGPT
\cite{lin2024videodirectorgpt} uses a large hosted LLM purely as a
planner, expanding one prompt into a structured video plan (per-scene
descriptions, entity layouts, cross-scene consistency groupings) that
drives a separate grounded video generator --- architecturally the closest
published analogue to this project's own Attempt 22 (a small local model
planning a multi-entity scene for a separate, deterministic renderer),
though built on models far larger than fit this project's card and
without this project's own compositional-generalization diagnosis. ViMax
\cite{vimax2025} is a broader, more recent agentic script-to-video
pipeline with no VRAM/scale figures found in the survey.

\paragraph{Small video diffusion backbones.} Wan2.1-T2V-1.3B
\cite{wanai2025wan21} is reported at roughly 8.19~GB inference VRAM at
FP16 (one source; another states 11~GB for the Diffusers pipeline
specifically), with GGUF-quantized builds down to 4--6~GB at 480p; LoRA
training, however, was only demonstrated on a 24~GB card in the sources
surveyed, and whether it fits under 8--12~GB is explicitly unverified.
CogVideoX-2B \cite{cogvideox2b} is reported at roughly 8~GB at FP8 for
480p/49-frame/8fps inference.

\paragraph{Layered/vector animation generation.} AniClipart
\cite{wu2024aniclipart} animates a static clipart image via B\'ezier
curves over keypoints with a text-to-video diffusion prior (Score
Distillation Sampling) as the motion signal and an explicit rigidity
constraint; Differentiable Motion Trajectory \cite{zhu2025dmt} does the
analogous thing for stroke-based sketches. Neither is shown handling
multiple independent characters plus a background --- the survey's own
reading is that multi-character vector-animation composition is an open
gap in the published literature, which is the gap Attempt 21 (\secsym9)
closes for this project's own material, by a route (deterministic
symbolic compositing, not diffusion or SDS) not covered in this survey.

\paragraph{Datasets.} Sakuga-42M \cite{sakuga42m2024} is the largest
hand-drawn animation corpus found (42M keyframes, $\sim$1.4M clips,
auto-generated captions, CC~BY-NC-SA~4.0). AnimeRun
\cite{animerun2022} and CreativeFlow+ \cite{creativeflowplus} supply dense
correspondence/optical-flow labels at far smaller scale. HumanML3D
\cite{humanml3d} and AMASS \cite{amass} supply captioned 3D human motion,
not 2D or stick-figure native; Motion-X/Motion-X++ \cite{motionx2025} is
reported larger than HumanML3D with subsets already reformatted to its
schema. This project's own corpus, LottieAnimation-660K
\cite{lottieanimation660k} (already CC-BY-NC-SA-4.0 licensed and in active
use), is confirmed by the survey as the right choice among these for
vector, captioned data with free shape correspondence, with the known
limitation (this project's own Attempt 16) that it is motion-graphics
iconography containing no characters.

\paragraph{Small text-to-motion models.} MoMask \cite{momask2024}
(44.85M params), Light-T2M \cite{lightt2m2024} (4.48M params), and
TriC-Motion \cite{tricmotion2026} (13.86M params) are two to three orders
of magnitude smaller than even this project's smallest trained language
model (Qwen3-0.6B), suggesting VRAM is unlikely to be the binding
constraint for this class of motion model; MotionGPT
\cite{motiongpt2023} is excluded from this size class because it is built
on an LLM backbone.

\paragraph{Evaluation.} FVD, and its content-bias correction
\cite{fvdcontentbias2024} (content-debiased FVD \cite{contentdebiasedfvd};
JEDi \cite{jedi2024}), are the standard distributional video-quality
metrics but require both a pretrained I3D-family feature network and, per
the survey's own reading, ``hundreds to thousands of clips per condition''
for a stable estimate --- a sample size this project's few dozen clips per
condition does not meet, and the reason FVD is not attempted here
(\secsym4). CLIPSIM (rooted in CLIPScore \cite{hessel2021clipscore}, first
applied to video in GODIVA \cite{godiva2021}) and its newer alternative
VQAScore/t2v\_metrics \cite{vqascore} measure text-frame alignment;
warping error (Farneback- or RAFT-based optical flow, warp-then-diff) is
the standard blind temporal-consistency metric and is exactly the quantity
this project's own Attempt 18 first approximated as a raw inter-frame-diff
ratio. Fr\'echet Video Motion Distance (FVMD) \cite{fvmd} isolates motion
consistency specifically; EvalCrafter \cite{evalcrafter2023} is a broader
benchmarking suite correlated against human judgment. This project's
evaluation harness (\secsym4) is built directly against this survey rather
than in ignorance of it, and documents where it diverges and why.

\paragraph{Script-to-multi-character animation at small scale (RESEARCH2
\secsym5).} This project's own second survey searched specifically for a
published system doing prompt-to-multi-character-animation at this
project's scale (sub-100M parameters, CPU-trainable) and found none:
AniMaker \cite{animaker2025} and AniME \cite{anime2025}, the most current
systems found in this space, are both multi-agent pipelines built on large
hosted LLMs and large video-diffusion backbones, not the ``one small
model, one CPU forward pass'' class this project occupies; the 2026
commercial storyboard-generator landscape surveyed (M~Studio, DomoAI,
TapVid, AnimateAI, Novi~AI) is entirely large-hosted-model product
surfaces. \textbf{As far as this survey found}, the prompt~$\rightarrow$~symbolic
scene script~$\rightarrow$~deterministic renderer pipeline this project
has built, entirely on CPU, is the smallest-scale system in this specific
niche in the literature to date. We state this as the survey states it,
hedge intact --- absence of a counter-example in one search pass is
evidence, not proof.

\paragraph{Compositional generalization and count extrapolation
(RESEARCH2 \secsym3).} The published record on fixing systematic
generalization failures like this project's cast-size wall (\secsym10.4--10.5,
\secsym11.10) is dominated by data-level, not architectural,
interventions: GECA \cite{andreas2020geca} took a seq2seq model on SCAN's
hardest ``jump'' split from unable to generalize at all to succeeding most
of the time via rule-based recombination of fragments seen in similar
contexts, and COGS/CFQ remained largely unsolved by naive scaling for
years afterward per the surveys RESEARCH2 checked. ``When Can Transformers
Count to $n$?'' \cite{yehudai2024counttransformers} gives a falsifiable,
checkable prediction this project used directly (\secsym11.4): exact
counting is reliably learnable when embedding dimension exceeds the
relevant vocabulary size, so a \texttt{d\_model} of 512 against closed
vocabularies of at most ten words rules out an embedding-width bottleneck
as the cause of this project's own counting/extrapolation limits. A
separate memory-augmented-counting source RESEARCH2 checked (exact
citation unresolved in that survey, reported there as unverified) found
plain transformer baselines generally ``incapable of extrapolation beyond
the training range'' on counting tasks --- independent corroboration that
this project's own cast-size wall is, per RESEARCH2's own framing, ``a
documented, general property of plain transformer/seq2seq architectures
$\ldots$ not something specific to t5-small or to this project's data,''
and we report it as confirming a known limitation rather than as a novel
defect (\secsym11.10).

\paragraph{Constrained decoding (RESEARCH2 \secsym4).}
Grammar/automata-guided constrained decoding is a well-studied, sound
technique class --- ABS \cite{abs2025} generalizes constrained beam search
to any constraint compilable to a finite automaton, which the project's
own duplicate-colour constraint (state = subset of ten colours used so
far) is trivially an instance of; Outlines \cite{outlines} and Guidance
\cite{guidance} integrate as logits processors directly against
HuggingFace \texttt{transformers} models, this project's own stack, and
RESEARCH2 recommended hand-writing the one constraint needed first and
holding those libraries as the fallback once more constraints accumulate
--- the recommendation this project followed (\secsym11.4) and later found
reason to reconsider once sampling replaced beam search and began emitting
tokens outside the DSL vocabulary entirely (\secsym11.9). The DETR-style
Hungarian set-decoder \cite{carion2020detr} is, per RESEARCH2 \secsym4,
the one technique surveyed that structurally addresses this project's
staging and duplicate-colour problems together --- $N$ cast slots
predicted with mutual attention make a repeated colour visible to the
model during generation in a way strict left-to-right decoding cannot, and
a repulsion term over the predicted positions attacks staging in the same
head --- but it is the only surveyed item requiring genuine architecture
work (a new decoder head, a new order-invariant matching loss) against
this project's ``minutes of CPU, no architecture change'' retraining
budget, and it was deliberately not built in Attempt 23 for that reason
(\secsym11.4). We report it as future work in \secsym13, with the
reasoning intact rather than as an unexplained gap. (Related, but not
directly adopted: GBNF/llama.cpp grammars \cite{llamacppgbnf}; factor-graph
layout constraints \cite{factorgraphs2024}; contextualized layout
refinement \cite{zhao2020scenegraphlayout}; GDPP diversity loss
\cite{elfeki2019gdpp}; LayoutFormer++ \cite{jiang2023layoutformerpp}; BLT
\cite{kong2022blt}; Grammar-Aligned Decoding/ASAp \cite{park2024asap};
Position Coupling \cite{cho2024positioncoupling}; Plan-and-Write
\cite{planandwrite2025}; PositionID \cite{positionid2024}; manga
panel-layout work \cite{mangapanel2024,comicsurvey2024}, surveyed and
found not to address within-panel character blocking, per RESEARCH2
\secsym2.)

\section{System Description}
\label{sec:system}

\subsection{Four representational eras (\texttt{ATTEMPTS.md})}
\label{sec:eras}

\textbf{Era 1 --- equations.} \texttt{geovid.py} renders scenes described
as parametric/explicit/polar/field numpy equation layers. A
\texttt{t5-small} (60M param) model maps prompts to a compact DSL spec,
trained on a synthesized dataset (no labeled prompt$\rightarrow$scene data
exists) --- 100\% per-slot validation accuracy, but capped to the 9
archetypes its own grammar defines. A hand-scripted gait engine
(\texttt{stickman.py}) produces the best-looking output of the project at
that stage, entirely by hand, with no learned component --- stated in the
log as the project's own indictment of Era 1--2: the models never beat the
script. (This is the standing indictment \secsym11 answers.)

\textbf{Era 2 --- distillation.} Two small nets (\texttt{PoseNet}, 81k
params, $<0.5^\circ$/joint error; \texttt{ImgPoseNet}, 70k params,
$\sim6^\circ$/joint error) are trained to imitate the hand-written engine
as teacher. A distilled student cannot exceed its teacher by construction;
in the shipped web app, the neural path is opt-in
(\texttt{EQV\_NEURAL=1}) because the hand-scripted engine renders more
reliably.

\textbf{Era 3 --- real data.} YOLO11-pose over 5000 COCO val2017 images
supplies 5051 real pose labels, genuinely useful as \emph{perception} but
not generation --- it does not produce new animation. \texttt{RealMotionNet},
trained on six Muybridge motion studies (3746 frames), loops smoothly but
blurs from cross-clip averaging; \texttt{MotionAR} (GRU, autoregressive,
val loss 0.0006) is sharper but non-looping, and is the shipped default
for learned motion.

\textbf{Era 4 --- AniSVG.} After a measured failure of per-frame raster
tracing of public-domain cartoons (\secsym3.2), the project adopts a
symbolic vector format that declares a cast of shapes once per shot and
encodes every frame as integer-delta edits to that cast --- translate,
rotate, scale, opacity, per-vertex morph, hide/show. Because shape $k$ is
the same shape at every timestep by construction, this format gets shape
correspondence for free where per-frame raster tracing could not, and is
what makes low-token-cost delta encoding possible at all. Era 4 is also
where the project's two arcs (\secsym1) diverge: AniSVG-trained language
models pursued styled character generation via a diffusion hybrid
(\secsym\secsym7--8); the scene-script system of Attempts 21--23
(\secsym\secsym9--11) instead composes the project's original Era~1 rig
(\texttt{stickman.py}) directly, without a learned appearance model, and
reaches multi-character animation first.

\subsection{Why raster tracing was rejected (Attempt 11)}
\label{sec:raster}

Public-domain Fleischer cartoon frames were traced to per-frame SVG with a
full built pipeline (crop/posterize/palette-fit/trace/rasterize/encode).
Measured on 144 real traced frames: raw SVG costs 6223 tok/frame; the best
compact encoding reaches 3794 tok/frame ($1.64\times$ compression). Delta
encoding against the previous frame bought almost nothing ($1.61\times
\rightarrow 1.64\times$) because independently-traced frames have no shape
correspondence to diff against --- frame-to-frame path correspondence in
raster-traced animation is an open research problem. A
background/foreground split made things \emph{worse} ($0.49\times$)
because these cartoons pan and truck constantly (35--80\% of pixels move
per shot) and the ragged motion-mask boundary adds more path complexity
than it saves. At $\sim$4000 tok/frame, 15~fps costs $\sim$60{,}000
tokens/second of video --- far beyond what a small model ($\sim$2000
coherent tokens generatable, i.e.\ $\sim$40--60 tok/frame for a short
clip) can produce. This measured failure is what motivated AniSVG: what
full-scene raster tracing cannot supply --- persistent shape identity
across frames --- is exactly what a declarative cast-of-shapes format
supplies for free.

\subsection{AniSVG format and corpus}
\label{sec:anisvg}

\begin{verbatim}
H <w> <h> <fps> <nframes>                 header
P <hex> ...                               palette, index = position
S <id> <pal> <sw> <x0> <y0> <dx dy>...    cast member; sw = stroke width x10
@ <frame> <op> ...                        per-frame edits; unlisted shapes hold
\end{verbatim}

The corpus source is \texttt{LottieGPT/LottieAnimation-660K} (660k
text-captioned Lottie/Bodymovin JSON files, CC-BY-NC-SA-4.0). Sampling a
Lottie shape tree at successive times gives shape correspondence for free
--- the property raster tracing could not supply. Only the metadata shard
(2.96~GB of the repository's 132~GB; the video-preview shard is never
touched) is needed. Filtering (cast size 2--20, no unsupported Lottie
feature, motion in at least half the frames, not static) retains 20\% of
a 1500-record pilot, projecting to roughly 130k usable clips from 660k.
Final corpus: \textbf{68/68 shards, 91{,}387 clips, 96~MB gzipped}; packed
to \textbf{78{,}971 train/val clips, 198.3M tokens}, with over-long clips
trimmed rather than dropped.

Measured cost on real corpus data (400 clips, Qwen3 tokenizer): median 71
tok/frame, 1768 tok/clip. Compared against Attempt 11's raster tracing:
$\sim19\times$ cheaper on a matched 36-shape, 48-frame Lottie ($\sim$4000
tok/frame $\rightarrow$ $\sim$213 tok/frame).

\subsection{The training bottleneck was VRAM from the loss, not the weights}
\label{sec:vrambottleneck}

\texttt{Qwen3ForCausalLM} upcasts logits to fp32 for cross-entropy against
a 151{,}936-token vocabulary; a single $\sim$1850-token sample allocates
$\sim$1.05~GB for the logit tensor alone, plus the same again for its
gradient. Scoring the sequence in chunks (\texttt{chunked\_ce}) cut peak
VRAM from 4.92~GB to 3.16~GB at 0.6B for a 7\% speed cost, with builtin,
manual, and chunked loss implementations agreeing to 0.8340 on a fixed
batch.

\begin{table}[h]
\centering
\caption{Peak VRAM and speed by model size and quantization.
(Source: \texttt{ATTEMPTS.md} Attempt 14 / \texttt{svg/README.md}.)}
\label{tab:vram-quant}
\begin{tabular}{@{}llrr@{}}
\toprule
Base & Quant & Peak VRAM & s/clip \\
\midrule
Qwen3-0.6B & bf16 & 3.16~GB & 0.76 \\
Qwen3-1.7B & bf16 & 5.30~GB & 1.32 \\
Qwen3-4B & NF4 & 6.28~GB & 4.00 \\
Qwen3-8B & NF4 + Soup layer streaming & $\sim$3.3~GB & see text \\
\bottomrule
\end{tabular}
\end{table}

Qwen3-8B was trained via \texttt{Soup} layer streaming, reported by its
authors at 119.6~tok/s on a 4~GB card, which puts a full 198.3M-token
epoch at roughly 19 days on this hardware --- the real run therefore
trained on a rate-measured subset rather than a full epoch.

\subsection{Two trained AniSVG generators, and what scale did not fix (Attempt 15)}
\label{sec:twogenerators}

\begin{table}[h]
\centering
\caption{v1 vs.\ v2 AniSVG generators. (Source: \texttt{ATTEMPTS.md} Attempt 15.)}
\label{tab:anisvg-v1v2}
\begin{tabular}{@{}lrr@{}}
\toprule
 & v1 Qwen3-8B (streamed, 1024 ctx) & v2 Qwen3-1.7B (resident, 4096 ctx) \\
\midrule
Clips seen & 28{,}034 & 77{,}397 \\
Final train loss & 0.4349 & \textbf{0.2709} \\
Valid generations & 94\% & \textbf{100\%} \\
Recognisable geometry & no & \textbf{still no} \\
\bottomrule
\end{tabular}
\end{table}

v2 saw roughly $7\times$ the animation signal of v1. The measured finding:
palette and op structure are learned; vertex sequences (long runs of exact
integers where one wrong digit deforms the shape) are not --- this is
\textbf{coordinate regression through next-token prediction}, not a
model-size or corpus-size problem, and it is what motivates splitting
motion (language model) from appearance (diffusion) in the arc-one hybrid
(\secsym3.7, \secsym\secsym7--8).

\subsection{A generated corpus, since no real anime vector-character corpus exists (Attempt 16)}
\label{sec:generatedcorpus}

LottieAnimation-660K is motion-graphics iconography with no characters; a
search for real anime vector-animation data at usable scale found none.
The project therefore generates its own character corpus from
\texttt{stickman.StickFigure}: \texttt{chars.py} (20{,}000 line-art
stick-figure clips) and \texttt{anime\_chars.py} (20{,}000 cel-style
clips: oversized head, hair silhouette, large low-set eyes, tapered
limbs). The cel figure's fixed 16-part schema in fixed paint order is
later what lets a pose skeleton be recovered from text a model wrote
(\secsym7.1).

\subsection{The diffusion hybrid: motion from a language model, style from diffusion (Attempt 17)}
\label{sec:hybrid}

The project's stated reasoning: the token budget is the ceiling --- a
small model affords $\sim$40--80 tokens/frame, motion is cheap in tokens
and appearance is not. The split: the language model generates motion
(AniSVG deltas), a diffusion model renders appearance, conditioned on
exact guides (\texttt{guides.py}) rather than a per-frame detector
estimate. Feasibility was measured: Wan~1.3B and AnimateDiff-Lightning run
at inference on 8~GB; LTX-2 LoRA \emph{training} needs 32--80~GB, ruling
out training a video backbone on this hardware while running a pretrained
one does not. Results and the eventual closure of this route are
\secsym\secsym7--8.

\section{Evaluation Harness (\texttt{critic/})}
\label{sec:critic}

The project built an eight-metric evaluation harness rather than trusting
loss curves or visual inspection, explicitly cross-referenced against
\texttt{paper/RESEARCH.md}'s survey of published metrics.

\begin{table}[h]
\centering
\small
\caption{This critic's metrics vs.\ the published metrics they stand in
for. (\secsym2)}
\label{tab:critic-vs-published}
\begin{tabular}{@{}L{3.3cm}L{4.2cm}L{6.5cm}@{}}
\toprule
This critic & Published metric & Why substituted \\
\midrule
\texttt{temporal\_stability.warp\_error} & warping error (optical-flow
  warp-then-diff) & same definition, Farneback flow (CPU) instead of RAFT
  (GPU) --- an extra deep net for the judge was ruled out on this
  hardware \\
\texttt{temporal\_stability.flicker\_ratio} & the raw inter-frame-diff
  ratio already computed in Attempt 18 & reproduced verbatim as a metric
  rather than reinvented \\
--- & FVD / content-debiased FVD / JEDi & needs a pretrained I3D-family
  net and ``hundreds to thousands of clips per condition'' for a stable
  estimate; this project has at most a few dozen per condition \\
--- & CLIPSIM / VQAScore & implemented (metric 6) but shown to be a null
  result in this domain (\secsym4.3) \\
--- & FVMD & same sample-size problem as FVD \\
\texttt{motion\_presence}, \texttt{identity\_preservation},
  \texttt{pose\_fidelity}, \texttt{scene\_complexity},
  \texttt{scene\_presence\_audit} & not in the published set & added
  because none of FVD/CLIPSIM/warping-error catch this project's specific
  known failure modes (a frozen clip scoring perfectly, an identity that
  reorganizes mid-clip, a render that ignores its pose guide, a
  multi-character scene a photo-trained detector cannot see at all) \\
\bottomrule
\end{tabular}
\end{table}

Metrics 1--5 (temporal stability, identity preservation via HSV histogram
cosine similarity, motion presence via Farneback flow coverage/energy,
pose fidelity via YOLO11-pose PCK, scene complexity via person-count
stability) each carry explicit pass/hard-fail thresholds justified in
\texttt{critic/README.md}. Metric 8, \texttt{scene\_presence\_audit}
(\secsym4.6), was added in response to Attempt 21's own finding that
\texttt{scene\_complexity} cannot see this project's material at all. We
restate here only the results that bear on the paper's central claims, and
flag every number that has been corrected since an earlier evaluation
pass.

\subsection{Motion-floor calibration (sanity check)}
\label{sec:motionfloor}

A synthetic clip built by repeating one real frame $30\times$ verifies the
harness does not reward a frozen output: \texttt{warp\_error} 0.00001,
identity 1.00, but \texttt{motion\_coverage} 0.0025\% --- correctly gated
\texttt{NOT-AN-ANIMATION: \ldots\ stable because nothing moved} against a
0.20\% floor. (Source: \texttt{critic/README.md} \secsym``Motion
presence''.)

\subsection{The baseline table (22 clips, current)}
\label{sec:baselinetable}

The following is reproduced verbatim from \texttt{critic/README.md}'s
current 22-clip baseline table --- \textbf{every \texttt{id worst} value
below is post-fix} for the crop-regime bug described in \secsym5.2, and
where a prior version of this table (an earlier evaluation pass) reported
a different number, that earlier number is superseded and not used
anywhere else in this paper.

\begin{table}[h]
\centering
\footnotesize
\caption{Baseline table (selected rows, of 22 total; full table in
\texttt{critic/README.md}). (Source: \texttt{critic/README.md}, ``Baseline
table.'')}
\label{tab:baseline}
\begin{tabular}{@{}L{4.3cm}rrrrrl@{}}
\toprule
Clip & $n$ & warp\_err. & id worst (regime) & mot.\ cov. & CLIPSIM & Verdict \\
\midrule
\texttt{control\_wave.mp4} (control) & 30 & 0.0011 & 1.00 (fallback) & 1.01\% & 0.329 & USABLE \\
\texttt{control\_kick.mp4} (control) & 30 & 0.0048 & 1.00 (fallback) & 3.29\% & 0.233 & USABLE \\
\texttt{out/a21\_park.mp4} (A21, 3 char.) & 200 & 0.0055 & 1.00 (fallback) & 5.32\% & 0.311 & \textbf{USABLE} \\
\texttt{out/a21\_street.mp4} (A21, 4 char.) & 200 & 0.0065 & \textbf{1.00} (fallback) & 4.89\% & 0.336 & \textbf{USABLE} $\leftarrow$ flipped, was PROMISING at 0.4368 \\
\texttt{corpus-anime.mp4} (reel) & 40 & 0.0711 & \textbf{1.00} (fallback) & 26.9\% & N/A & \textbf{PROMISING} $\leftarrow$ flipped, was NOT-AN-ANIMATION at 0.29 \\
ground truth \texttt{10887279} & 18 & 0.0088 & 0.84 (no YOLO) & 10.4\% & 0.337 & USABLE \\
\texttt{svg/out/gen17/g002} & 28 & 0.0089 & 0.99 (fallback) & 21.1\% & 0.306 & \textbf{NOT-AN-ANIMATION} $\leftarrow$ flipped by structural metric \\
\texttt{svg/out/styled/frames} (A18 hybrid) & 8 & \textbf{0.1408} & 0.63 (fallback) & 60.4\% & 0.283 & \textbf{NOT-AN-ANIMATION} \\
\texttt{critic/tmp/static\_test.mp4} (neg.\ control) & 30 & 0.00001 & 1.00 (fallback) & 0.003\% & N/A & \textbf{NOT-AN-ANIMATION} \\
\bottomrule
\end{tabular}
\end{table}

\subsection{CLIPSIM cannot separate content-correct from content-wrong generations (null result)}
\label{sec:clipsim}

The calibration test: score the hand-written positive controls
(known-correct content) and a known-bad generated clip (\texttt{g002},
captioned ``two four-pointed star-shaped figures\ldots'', which draws two
blobs --- neither a star, per Attempt 15) with CLIP ViT-B/32
(\texttt{openai/clip-vit-base-patch32}, CPU-only), against their true
captions.

\begin{table}[h]
\centering
\caption{CLIPSIM calibration. (Source: \texttt{critic/README.md}
\secsym``Content correctness / prompt alignment (CLIPSIM)''.)}
\label{tab:clipsim}
\begin{tabular}{@{}L{5.5cm}L{3.5cm}r@{}}
\toprule
Clip & Caption used & CLIPSIM mean \\
\midrule
\texttt{control\_wave.mp4} (correct content) & its own correct caption & \textbf{0.329} \\
\texttt{control\_kick.mp4} (correct content) & its own correct caption & \textbf{0.233} \\
\texttt{g002} (known-bad: blobs, not stars) & its own correct caption & \textbf{0.306} \\
\texttt{control\_kick.mp4} & wrong caption (wave's) & 0.235 \\
\texttt{control\_wave.mp4} & unrelated caption (``a car\ldots'') & 0.163 \\
\bottomrule
\end{tabular}
\end{table}

The known-bad blob clip scores \emph{higher} (0.306) than the known-good
kick control (0.233): no single threshold passes both controls and fails
the blobs. \texttt{CONTENT\_CORRECTNESS\_FLOOR = 0.18} is kept only as a
``completely different concept'' tripwire, not as the content-correctness
check it was built to be; over the full baseline table, adding CLIPSIM
changed \textbf{zero} verdicts.

\subsection{Structural correctness separates the same pair by \texorpdfstring{$\sim27\times$}{~27x} (positive result)}
\label{sec:structural}

Because a generated AniSVG clip's output \emph{is} its source --- a
declared cast of shapes in a fixed paint order --- content correctness can,
on this one path, be checked exactly. The check (\texttt{star\_score})
measures the discrete Fourier coefficient of a shape's
centroid-normalized radius sequence at the frequency the caption implies.

\begin{table}[h]
\centering
\caption{Structural correctness calibration. (Source: \texttt{critic/README.md}
\secsym``Structural correctness (metric 7)''.)}
\label{tab:structural}
\begin{tabular}{@{}L{4.5cm}crc@{}}
\toprule
Clip & Shapes $\geq0.05$ (star-like) & Shape count & \texttt{geometry\_verified} \\
\midrule
ground truth (\texttt{10887279}) & 3, at 0.135/0.154/0.140 & 11 & \textbf{True} \\
\texttt{g002} (blobs) & 0, all four at 0.0048--0.0052 & 4 & \textbf{False} \\
\bottomrule
\end{tabular}
\end{table}

The weakest genuine star shape (0.135) and the strongest blob shape
(0.0052) differ by roughly $\mathbf{27\times}$, with no tuning pressure on
the 0.05 threshold. Adding this metric changes exactly one verdict in the
table: \texttt{g002} flips from USABLE (passing five pixel metrics and
CLIPSIM) to \textbf{NOT-AN-ANIMATION}, correctly.

\subsection{The evaluation asymmetry, stated explicitly}
\label{sec:asymmetry}

Metrics 6 and 7, read together, are one of the harness's central
findings: one is a null result (CLIP embedding similarity cannot separate
the populations in this domain); the other separates the same pair
cleanly, at zero GPU cost, because it parses declared source rather than
inferring content from pixels. This is \textbf{not a general solution to
content correctness} --- it is specific to the one output path where
generation and source are the same text. The diffusion/hybrid path
(\texttt{svg/out/styled}) has no declarative source to check against; it
is rejected on temporal grounds with total confidence, but this harness
would have no way to verify a future, temporally-fixed version of that
pipeline actually painted the character its prompt asked for. Attempt
21's own composed scenes carry the same asymmetry in a new form
(\secsym9.7): a scene script is itself a declarative source and could in
principle be checked exactly the same way AniSVG is, but nothing does
that yet.

\subsection{A detector-free presence audit, added because \texttt{scene\_complexity} cannot see this project's own material (metric 8)}
\label{sec:presenceaudit}

\texttt{yolo11n-pose.pt}, the person detector \texttt{scene\_complexity}
(metric 5) is built on, finds \textbf{no person at all} in either of this
project's own hand-written single-figure controls (\texttt{control\_wave},
\texttt{control\_kick}, \texttt{modal\_count: 0} on both) and the same is
true of hand-scripted scenes containing three or four visible,
correctly-drawn characters (\texttt{out/a21\_park.mp4},
\texttt{out/a21\_street.mp4}). (Source: \texttt{critic/README.md}
\secsym``Scene complexity''.) A photo-trained detector does not read flat
stick-figure line art, at any cast size --- reporting \texttt{modal\_count:
0} in that situation is indistinguishable from ``this clip genuinely has
no characters,'' which is false every time it has happened in this
project. \texttt{scene\_complexity} now applies the same
\texttt{MEASURABLE\_DETECTION\_RATE = 0.50} bar \texttt{pose\_fidelity}
already used, and reports \texttt{UNMEASURABLE} with the detection rate
attached rather than a false cast count of zero, below that bar; it
remains, as before, never hard-gated.

The working replacement, \texttt{scene\_presence\_audit} (metric 8),
verifies presence rather than detecting it: given the declared colour of
each cast member (supplied by the caller, not inferred), it checks colour
distance under a tolerance and connected-component count per frame --- the
same technique \texttt{scenescript.py --audit} (\secsym9) uses,
reimplemented independently in \texttt{critic/} so the two tools stay
comparable without a code dependency.

\begin{table}[h]
\centering
\caption{Scene presence audit. (Source: \texttt{critic/README.md}
\secsym``Scene presence audit''.)}
\label{tab:presenceaudit}
\begin{tabular}{@{}lcL{4cm}L{4cm}@{}}
\toprule
Clip & Characters & All visible (this tool, 90-frame sample) & All visible
  (\texttt{scenescript.py --audit}, full clip) \\
\midrule
\texttt{a21\_park.mp4} & 3 & 66/90 (73\%) & 75/75 sampled (100\%) \\
\texttt{a21\_street.mp4} & 4 & 61/90 (68\%) & 84/84 sampled (100\%) \\
\bottomrule
\end{tabular}
\end{table}

The two tools disagree on the exact rate because they sample differently,
not because either mistracks presence; both agree that every declared
character is visible in the large majority of sampled frames --- the
thing \texttt{scene\_complexity} cannot see on this art at all. This
metric is, explicitly, a verification tool rather than a detector: it has
nothing to measure without being told what to look for, so it cannot
replace \texttt{scene\_complexity} on arbitrary or unlabelled input, and
it is never gated.

\section{Evaluation Methodology as a Contribution}
\label{sec:evalmethodology}

This project's evaluation harness and decoding pipeline produced four
distinct false or unmeasurable verdicts in the course of the work. We
report all four as a contribution in their own right, because the
underlying failure mode --- a metric trained or calibrated on
photographic content, or a decoding default tuned for fluent free text,
silently mismeasuring or misrepresenting stylized, symbolic, procedurally
generated, or unsupervised-slot content --- is not specific to this
project's tooling, and we believe it generalizes to anyone evaluating
generative 2D animation or any generative model with unsupervised output
slots.

\subsection{Case 1 --- CLIPSIM cannot separate correct from incorrect content in this domain (\texorpdfstring{\secsym4.3}{4.3})}
\label{sec:case1}

Already described above: a known-wrong ``blob'' generation (0.306) scores
\emph{higher} than a correctly-drawn hand-written control (0.233), and
giving a correct clip the wrong caption barely moves its score ($0.233
\rightarrow 0.235$). CLIP ViT-B/32 at CPU budget, with no domain
fine-tuning, provably cannot distinguish ``kicking'' from ``waving,'' let
alone a correct star from a blob, in this project's stylized 2D domain.
This is reported as a null result rather than silently dropped or its
threshold tuned until it looked like it worked.

\subsection{Case 2 --- a crop-regime artefact that cost one real clip its verdict twice}
\label{sec:case2}

\texttt{identity\_preservation} (metric 2) crops ``the character'' per
frame --- a YOLO detection box when one fires, a 70\% centre-crop fallback
otherwise --- and previously chose between the two \textbf{independently,
frame by frame}. On \texttt{out/a21\_street.mp4} (four hand-scripted
characters, Attempt 21), YOLO fired on exactly 1 of 24 sampled frames (a
$48\times135$ box around one small figure); the other 23 frames used the
identical centre-crop box. The metric then compared that one
geometrically unrelated crop against the rest and reported
\texttt{worst\_similarity: 0.4368}, under the 0.45 pass bar, capping the
clip at PROMISING. \textbf{The clip never changed. Only the measurement
basis changed, once, mid-clip}, and the metric reported that as an
identity change; an independent colour-based audit of the same clip found
all 4 characters visible in 84 of 84 sampled frames. Forcing the fallback
crop for every frame gives \textbf{0.9972} on the same clip. (Source:
\texttt{critic/README.md} \secsym``Bug found and fixed: mixed crop
regimes''.)

\textbf{The fix} (\texttt{character\_boxes()}) chooses exactly one crop
regime for the whole clip and never mixes them, using the same
\texttt{MEASURABLE\_DETECTION\_RATE = 0.50} / \texttt{DETECTION\_CONF\_FLOOR
= 0.40} bar \texttt{pose\_fidelity} already applied: below 50\% confident
detection, every frame uses the fallback crop; at or above it, every box
is YOLO-derived, with gaps held forward or backward in time rather than
falling through to the fallback. The regime used, why, and the
confident-detection rate behind that choice are always reported
(\texttt{crop\_regime}, \texttt{crop\_regime\_reason},
\texttt{confident\_detection\_rate}).

\textbf{A second, previously undiagnosed instance of the same bug was
caught retroactively once the fix was applied.}
\texttt{release/corpus-anime/corpus-anime.mp4} had been scoring
\texttt{NOT-AN-ANIMATION} on
\texttt{identity\_preservation.worst\_similarity: 0.29}, attributed at the
time to genuine identity drift across a multi-clip reel --- a wrong
explanation that had been carried in this project's own documentation with
confidence. Re-run with the fix, the same clip scores
\texttt{worst\_similarity: 0.99995}, \texttt{crop\_regime: fallback}
throughout (confident detection rate 20\%, below the 50\% bar); the 0.29
was the same crop-mixing artefact, not identity drift of any kind, and the
clip is now \texttt{PROMISING}, capped only by \texttt{warp\_error}
(0.0711 against a 0.06 bar), which was never in dispute. (Source:
\texttt{critic/README.md}, retraction under ``Most damning finding.'') We
state the general lesson explicitly, as the project's own documentation
now does: \textbf{the failure signature of this bug class is a
near-perfect score on almost a whole clip dragged through a pass bar by
one or a few outlier samples measured a different way} --- if removing a
single sample from a ``worst-case'' statistic changes the verdict, and
that sample was measured on a different basis than its neighbours, the
metric, not the clip, is what changed.

\subsection{Case 3 --- person/pose detectors read zero on this project's stylized art, independent of what is actually there}
\label{sec:case3}

CMU OpenPose (\texttt{controlnet\_aux}) returns no person at all on this
project's rendered skeletons, flat cel renders, or SD1.5 anime line-art
output, while reading a real photograph fine (12/12 keypoints on
\texttt{zidane.jpg}); YOLO11-pose does somewhat better but detected a cel
figure confidently in only 2 of 6 tested clips (\secsym8.3, Attempt 19).
The same detector finds \textbf{no person in any frame} of
\texttt{out/a21\_park.mp4}, a clip that genuinely contains three
characters, and one person in 1 of 200 frames of
\texttt{out/a21\_street.mp4}, which contains four --- and, checked
directly as a control before drawing any conclusion, the identical
\texttt{modal\_count: 0} result also appears on this project's own
single-figure hand-written positive controls. (Source:
\texttt{ATTEMPTS.md} Attempt 21, Result 3; \texttt{critic/README.md}
\secsym``Scene complexity'' correction.) The root cause, established
independently in Attempt 19 (\secsym7.1), is the same proportion mismatch
that makes ControlNet ignore this project's pose guides: every detector
available here was trained on photographs of human proportions, and this
project's line-art rig sits outside that distribution regardless of how
many correctly-drawn characters are on screen. \texttt{scene\_complexity}
was corrected to report \texttt{UNMEASURABLE} with the detection rate
attached rather than a false cast count of zero (\secsym4.6).

\subsection{Case 4 --- beam search collapsing an adequate learned distribution into a false ``the model cannot stage'' verdict}
\label{sec:case4}

Every position number reported in Attempt 22, and initially in Attempt
23's own ablation cells, was decoded with \texttt{num\_beams=4}. Beam
search returns the approximately most-likely \emph{sequence}; for an
output slot the prompt does not constrain (character $x$-position, in this
project's DSL), the most likely single value is, by definition, the
marginal mode of whatever the model actually learned --- so ``the model
emits $x\approx9.4$ nearly every time'' is consistent with two different
situations that call for opposite fixes: the model's conditional
distribution has genuinely collapsed (a data or architecture problem), or
the distribution is fine and the search is standing on its peak (a
one-argument decoding fix). Nothing in Attempt 22 distinguished these.
(Source: \texttt{ATTEMPTS.md} Attempt 23 \secsym6.)

Decoding the \emph{same checkpoint} on the \emph{same 30 held-out
prompts} four ways and scoring each with an independently-calibrated
staging metric (\secsym11.2) answered it directly:

\begin{table}[h]
\centering
\caption{Decode probe: staging and distinct $x$-values by decoding
strategy, on identical weights and prompts.}
\label{tab:decodeprobe}
\begin{tabular}{@{}lrrrr@{}}
\toprule
Decode & v3 staging & v3 distinct $x$ & v2 staging & v2 distinct $x$ \\
\midrule
\texttt{beam4} (prior default) & \textbf{0.188} & \textbf{3} & \textbf{0.333} & 6 \\
\texttt{greedy} & 0.189 & 3 & 0.393 & 7 \\
\textbf{\texttt{sample}} (top-$p$ 0.95, $T{=}1.0$) & \textbf{0.862} & \textbf{55} & \textbf{0.845} & 20 \\
\texttt{sample\_lo} (top-$p$ 0.90, $T{=}0.7$) & 0.732 & 43 & 0.796 & 21 \\
\bottomrule
\end{tabular}
\end{table}

Switching only the decoding argument --- same weights, same prompts ---
moves staging from 0.188 to 0.862 on one checkpoint and 0.333 to 0.845 on
another, both landing \emph{above} the weaker of the project's two
hand-written reference clips (0.825). Thirty scenes decoded by beam search
produce three distinct $x$-coordinates in total across the whole set, with
87.8\% of all characters landing on one of them; the identical weights,
sampled, produce fifty-five. \textbf{A model that had genuinely collapsed
could not produce fifty-five distinct positions; a search that maximizes
sequence likelihood produces three whether or not the model had
collapsed.} This corrects a specific, previously-published claim in this
project's own record (Attempt 22 \secsym10.1: staging is ``the single
biggest defect in every generated scene $\ldots$ the decoder emits modal
values and characters cluster,'' attributed to absent supervision) --- the
absent supervision is real, but the \emph{reported symptom} was mostly an
artefact of how the earlier numbers were decoded, not of what the model
had learned.

\textbf{A companion ablation shows what would have been reported, and
believed, without this check.} A deterministic post-decode spacing pass
--- plausible engineering, and RESEARCH2's own top-ranked recommendation
for this defect --- moves beam-search staging from 0.178 to 0.587 and
looks, in isolation, like a successful fix for a model limitation. Reading
it alongside the decoding-only result shows what it actually is:

\begin{table}[h]
\centering
\caption{The honest decomposition of the staging number (in
distribution).}
\label{tab:honestdecomp}
\begin{tabular}{@{}L{7.5cm}r@{}}
\toprule
Source of the staging number & Staging (in dist.) \\
\midrule
Model output, read by beam search & 0.178 \\
$\ldots$ plus a deterministic spacing pass enforcing it after decoding & 0.587 \\
\textbf{Model output, read by sampling instead --- no pass at all} & \textbf{0.862} \\
Hand-written reference clips & 0.825--1.000 \\
\bottomrule
\end{tabular}
\end{table}

Under beam search, the deterministic pass does essentially all of the
visible work and never reaches the hand-written band (only 17\% of
scenes clear the worse reference clip); changing the decoder alone
reaches 0.862, past that same clip, because it is the model's own answer
rather than an imposed one. \textbf{We state this exactly as the project
states it: good engineering, bad science claim.} Without the decoding
ablation, this project would have shipped the spacing pass and reported
that the model had learned to stage a scene --- a claim the same
project's own later measurement shows to be false. (Source:
\texttt{ATTEMPTS.md} Attempt 23 \secsym7, \secsym11.7.)

\subsection{The general lesson}
\label{sec:generallesson}

Read together, these four cases motivate a single methodological point we
generalize beyond this project: \textbf{a metric's or a decoding
default's provenance --- what it was trained, calibrated, or tuned for
--- bounds what it can be trusted to measure or produce}, and this bound
does not announce itself; it must be checked against a known-good and a
known-bad case, or a same-weights control decoded a second way, before
the result is trusted as a verdict on the underlying system. CLIP ViT-B/32
was never going to separate stars from blobs in flat vector art without
that check; a YOLO person detector was never going to count stick
figures; a metric that silently switches its own measurement basis
mid-series will produce a result that looks like a property of the clip
and is actually a property of the metric; and a decoding strategy that
reports only the mode of a learned distribution will produce a result
that looks like a property of the model and is actually a property of the
search. All four were caught in this project only because a control --- a
known-good clip, a hand-authored ablation, a re-run after a fix, or the
same weights decoded a second way --- was run and compared, never
assumed. We consider this pattern of self-correction (\secsym1; also
Attempt 19's topology-to-proportion correction, \secsym7.1) as important a
contribution as any single result it produced.

\section{Reproducibility}
\label{sec:repro}

\textbf{Hardware.} RTX 4060 Laptop GPU, 8~GB VRAM, run under a documented
50\%-duty-cycle GPU guard (\texttt{tools.gpuguard.Guard}) for thermal
reasons --- the laptop's fan is damaged and an elevated-shell clock lock
was unavailable, so duty cycling carried thermal management alone for
every GPU loop reported in this paper. \texttt{USE\_TF=0} is required
project-wide or transformers 4.57 attempts a Keras-3 backend and fails.
\textbf{Attempts 21--23 use no GPU at any point except the critic's own
YOLO pass} --- scene composition and every version of the script-writing
model are CPU-only, and this is reported as a property of the route, not
an incidental detail (\secsym9, \secsym10, \secsym11).

\textbf{Model versioning.} Three checkpoints are preserved under
\texttt{model/checkpoints/\{v1-confounded,v2-deconfounded,v3-staged\}/},
each with its own weights, the exact training data that produced it, and
a run summary, recorded in \texttt{model/MODELS.md}. Nothing in Attempt
23 writes into the v1/v2 directories; v3 trains to a new directory.
\texttt{model/verify\_checkpoints.py} proves, rather than asserts, that
the preserved copies still work: \textbf{v2-deconfounded reproduces
Attempt 22's published headline output byte-for-byte} on the prompt
``two friends meet in the park, one waves, then they kick a ball
around'' (\texttt{bg park dur 91 | cast ana blue 96 100 ; bo amber 92 90
| prop ball 58 100 ; tree 69 110 | tl 1 26 bo walk 92 ; 8 28 ana wave ;
28 47 bo kick ball ana ; 57 79 bo walk 92}, zero problems, zero repairs,
matching Attempt 22 \secsym8 exactly), and v1-confounded reproduces its
own documented failure mode (a cast of one referencing an undeclared
second character) on the same prompt --- a load producing anything else
would have been the alarm. (Source: \texttt{ATTEMPTS.md} Attempt 23
\secsym1.)

\textbf{Key commands} (all \texttt{USE\_TF=0}; see \texttt{svg/README.md},
\texttt{critic/README.md}, and \texttt{ATTEMPTS.md} Attempts 19--23 for
the complete set):

\begin{verbatim}
# corpus build and training (AniSVG / diffusion hybrid, arc one)
python svg/stream_corpus.py -o svg/data/corpus --max-clips 40000 --progress 2000
python svg/pack_dataset.py --corpus svg/data/corpus -o svg/data/train
python svg/train_lora.py --base Qwen/Qwen3-1.7B-Base --data svg/data/train
python svg/coco18.py
python svg/render_style.py --clip release/corpus-anime/clip00.anisvg \
    --guide coco18 --frames 8 -o svg/out/a19/styled-coco18-final

# AnimateDiff temporal module (Attempt 20)
python critic/evaluate.py svg/out/a20/w512_f8 --guide svg/out/a20/guides_w512/coco18 \
    --max-frames 8 -o svg/out/a20/critic/w512_f8
python svg/animate.py --sweep           -o svg/out/a20c-noofl
python svg/animate.py --sweep --offload -o svg/out/a20c-ofl
python svg/animate.py --frames 8 --width 512 -o svg/out/a20
python svg/animate.py --frames 8 --width 512 --ip-adapter --ip-scale 0.6 -o svg/out/a20-ip

# multi-character scene composition (Attempt 21) -- no GPU
python scenescript.py --selftest
python scenescript.py scenes/park_meet.scene    -o out/a21_park.mp4   --frames out/a21_park_frames
python scenescript.py scenes/street_relay.scene -o out/a21_street.mp4 --frames out/a21_street_frames
python critic/evaluate.py out/a21_park.mp4 --max-frames 200 --identity-max 24 \
    --caption "three stick figures in a park, two of them passing a ball" \
    --character-colors "50F0F0,FF963C,FF6EB4" -o critic/out/a21_park

# the model that writes scene scripts (Attempt 22) -- no GPU except the final critic pass
python model/scene_synth.py --n 6000 --val 700 --outdir model/data --prefix scene
python model/scene_synth.py --n 6000 --val 700 --outdir model/data --prefix scene2 --deconfound
python model/train.py --train model/data/scene2_train.jsonl --val model/data/scene2_val.jsonl \
    --out model/scene_ckpt2 --epochs 6 --batch 16 --lr 3e-4 --max-in 104 --max-out 176 --val-cap 200
python model/scene_infer.py "two friends meet in the park, one waves, then they kick a ball around" \
    --ckpt model/scene_ckpt2 -o scenes/a22_generated.scene
python scenescript.py scenes/a22_generated.scene -o out/a22_generated.mp4

# staging, decoding ablation, constrained decoding, and the two deliverables (Attempt 23)
python model/scene_synth.py --n 6000 --val 700 --outdir model/data --prefix scene3 --deconfound --stage
python model/train.py --train model/data/scene3_train.jsonl --val model/data/scene3_val.jsonl \
    --out model/scene_ckpt3 --epochs 6 --batch 16 --lr 3e-4 --max-in 160 --max-out 256 --val-cap 200
python model/scene_staging.py scenes/park_meet.scene scenes/street_relay.scene \
    --dir scenes --glob "a22_*.scene"
python model/scene_decode_probe.py --ckpt model/scene_ckpt3 --val model/data/scene3_val.jsonl --n 30
python model/scene_eval.py --ckpt model/scene_ckpt3 --val model/data/scene3_val.jsonl --n 150 \
    --out model/eval_out3 --max-in 160 --max-out 256 --sample --unique-colours --rerank 4 --tag _v3sr
python model/scene_probe.py --ckpt model/scene_ckpt3
python model/scene_layout_baseline.py \
    "model/eval_out3/in_distribution_v3.json|model/data/scene3_val.jsonl|v3"
python model/scene_infer.py \
    "two friends meet in the park, one waves, then they kick a ball around" \
    --ckpt model/scene_ckpt3 --max-in 160 --max-out 256 \
    --sample --unique-colours --rerank 4 --seed 23 -o scenes/a23_generated.scene
python scenescript.py scenes/a23_generated.scene -o out/a23_generated.mp4
python critic/evaluate.py out/a23_generated.mp4 --max-frames 200 --identity-max 24 \
    -o critic/out/a23_generated \
    --caption "two stick figures in a park, one waves, then they kick a ball to each other"
python model/scene_infer.py \
    "six friends spread out on a night street, all of them dancing, then two of them kick a ball around" \
    --ckpt model/scene_ckpt3 --max-in 160 --max-out 320 \
    --sample --unique-colours --rerank 4 --seed 5 -o scenes/a23_six.scene
\end{verbatim}

\textbf{Frame/artefact locations.} Attempt 20:
\texttt{svg/out/a20/w512\_f8/} (headline run), \texttt{svg/out/a20-ip/}
(+IP-Adapter), \texttt{svg/out/a20b/} (16-frame sweep),
\texttt{svg/out/a20c-noofl/} / \texttt{svg/out/a20c-ofl/} (VRAM
re-measurement). Attempt 21: \texttt{out/a21\_park.mp4},
\texttt{out/a21\_street.mp4}, contact sheets \texttt{out/a21\_sheet.png} /
\texttt{out/a21\_sheet\_street.png}. Attempt 22:
\texttt{model/eval\_out/} and \texttt{model/eval\_out2/} (every generated
string, parse error and repair, per prompt),
\texttt{model/eval\_out/clause\_probe.json} /
\texttt{clause\_probe2.json} (the confound probe),
\texttt{scenes/a22\_generated*.scene}, \texttt{out/a22\_generated*.mp4},
\texttt{critic/out/a22\_*}. Attempt 23: \texttt{model/MODELS.md} and
\texttt{model/checkpoints/\{v1-confounded,v2-deconfounded,v3-staged\}/};
\texttt{model/eval\_out3/} (every cell, \texttt{compare\_indist.txt},
\texttt{compare\_ood.txt}, \texttt{decode\_probe\_v\{2,3\}.json},
\texttt{clause\_probe3.txt}, \texttt{layout\_baseline.txt},
\texttt{staging\_calibration.\{json,txt\}}, \texttt{data\_report.txt});
\texttt{model/scene\_train3.log}; \texttt{scenes/a23\_generated.scene},
\texttt{scenes/a23\_six.scene}; \texttt{out/a23\_generated.mp4},
\texttt{out/a23\_six.mp4}; \texttt{critic/out/a23\_generated/},
\texttt{critic/out/a23\_six/}.

\textbf{A process failure worth recording as a reproducibility lesson in
its own right:} four Attempt 23 evaluation runs, written to the same
output directory under different flag combinations, silently overwrote
each other before the \texttt{--tag} argument existed on
\texttt{model/scene\_eval.py}; the loss was caught and the runs
re-executed with distinct tags. We record this rather than omit it
because it is the same class of failure \secsym5 documents for metrics
--- an unlabeled measurement basis silently replacing another --- applied
to raw experiment logging rather than to a metric's internal logic.

Every GPU loop in this paper ran through \texttt{tools.gpuguard.Guard} at
50\% duty; the clock lock needed an elevated shell and was unavailable, so
duty cycling carried thermal management alone throughout.
\texttt{USE\_TF=0} before every import; the existing interpreter was
reused across attempts, no new environment except the Python~3.12
\texttt{.venv-soup} for Soup layer streaming (\secsym3.4).

\section{Experiments and Results I --- The Diffusion Route, Closed by Measurement}
\label{sec:results1}

\subsection{A diagnosis corrected in the open: skeleton topology was not the cause}
\label{sec:diagnosiscorrected}

Attempt 18 rendered a procedural clip's exact pose guides through SD1.5 +
ControlNet-OpenPose (fixed seed, identical prompt, only the conditioning
image varying between frames):

\begin{table}[h]
\centering
\caption{Attempt 18 flicker ratio. (Source: \texttt{ATTEMPTS.md} Attempt 18.)}
\label{tab:a18flicker}
\begin{tabular}{@{}lr@{}}
\toprule
 & Inter-frame change \\
\midrule
Pose guides (input) & 1.29 \\
Styled frames (output) & 30.44 \\
\textbf{Flicker ratio} & $\mathbf{23.6\times}$ \\
\bottomrule
\end{tabular}
\end{table}

The style half worked --- genuine anime-style output at a quality no token
budget could buy directly. Two failures were logged: (1) the pose was
ignored; (2) identity was not preserved (hair, face, background reorganize
every frame despite a fixed seed). Attempt 18's own hypothesis for
failure (1): the hand-built skeleton was OpenPose-\emph{like} but not true
COCO-18 topology --- stated as ``fixable\ldots\ a mapping exercise.''

Attempt 19 did that mapping (\texttt{svg/coco18.py}) and verified it two
ways before spending any GPU time: the drawing constants were checked
against \texttt{controlnet\_aux.open\_pose.draw\_poses} directly, at
\textbf{max absolute pixel difference 0 over all 29 frames} of the test
clip; and the recovered joints matched the rig's own ground-truth joints
to sub-pixel accuracy (mean 0.40--0.62~px, max 1.10--1.28~px across four
actions, on a 256~px canvas).

\begin{table}[h]
\centering
\caption{Flicker ratio across four conditioning variants, no temporal
mechanism. (Source: \texttt{ATTEMPTS.md} Attempt 19.)}
\label{tab:a19flicker}
\begin{tabular}{@{}L{6.5cm}rrr@{}}
\toprule
Conditioning & Guide & Styled & Ratio \\
\midrule
A18: limb-axis skeleton, OpenPose palette & 1.29 & 30.44 & $\mathbf{23.6\times}$ \\
A19: true COCO-18, head from the art & 1.18 & 26.05 & $22.1\times$ \\
A19 + human-proportion head & 1.14 & 34.99 & $30.8\times$ \\
A19 + human head + per-clip side convention & 1.09 & 39.58 & $\mathbf{36.2\times}$ \\
\bottomrule
\end{tabular}
\end{table}

Row 2 is the direct test of Attempt 18's hypothesis: correct COCO-18
topology alone changed \textbf{nothing} --- the render still stands
hands-on-hips. \textbf{The topology diagnosis was right about what was
broken and wrong about what mattered.}

\textbf{The control that found the real cause.} Three conditioning
images through the identical pipeline:

\begin{table}[h]
\centering
\caption{The control ablation. (Source: \texttt{ATTEMPTS.md} Attempt 19.)}
\label{tab:controlablation}
\begin{tabular}{@{}L{7cm}l@{}}
\toprule
Conditioning & Result \\
\midrule
Genuine OpenPose annotation of a real photo (\texttt{zidane.jpg}) & obeyed \\
Hand-built canonical wave skeleton, arm straight up & \textbf{obeyed} \\
Rig-derived COCO-18, same wave pose & ignored \\
\bottomrule
\end{tabular}
\end{table}

A correctly-shaped COCO-18 skeleton \emph{is} obeyed by this model; the
fault remained in this project's own skeleton. Measured against the
obeyed reference, in units of torso length $T$:

\begin{table}[h]
\centering
\caption{Proportion mismatch, measured in units of torso length $T$.
(Source: \texttt{ATTEMPTS.md} Attempt 19.)}
\label{tab:proportions}
\begin{tabular}{@{}lrr@{}}
\toprule
Ratio & Ours (from the art) & Obeyed reference \\
\midrule
Neck $\rightarrow$ nose & 0.103~$T$ & 0.330~$T$ \\
Ear span & 0.956~$T$ & 0.400~$T$ \\
Shoulder width & 0.579~$T$ & 0.770~$T$ \\
\textbf{Ear span / shoulder width} & \textbf{1.65} & \textbf{0.52} \\
\bottomrule
\end{tabular}
\end{table}

\textbf{The real cause: anime proportions are not human proportions.} The
cel figure's head is deliberately drawn $1.65\times$ wider than its own
shoulders --- the style, by design. Mapped faithfully to canonical
topology, its metric proportions still sit outside anything a
human-photo-trained ControlNet has seen; the model discards the pose and
falls back on its prompt prior. We state this as a general finding
(contribution 7, \secsym1): the property that makes a character read as
anime is the same property that removes it from a human-pose conditioning
manifold. Two targeted corrections followed; obedience is still only
\textbf{partial} --- a raised arm survives clearly in one of four
inspected frames and is ambiguous in the rest. \textbf{Stated as unfixed,
not smoothed over.}

\subsection{Pose obedience and temporal stability are measured orthogonal, not correlated}
\label{sec:orthogonal}

The same four-row table above shows flicker rising \textbf{monotonically}
as pose conditioning improved: $23.6\times \rightarrow 22.1\times
\rightarrow 30.8\times \rightarrow 36.2\times$. Two contributing effects:
the canonical skeleton is drawn thinner than Attempt 18's, constraining
less of the frame; and pose conditioning is, by construction, orthogonal
to what drives flicker. \textbf{A better guide bought a strictly worse
flicker number, on a quantity the guide was never able to move.}

\subsection{An off-the-shelf pose estimator cannot score this material}
\label{sec:poseestimator}

CMU OpenPose returns \textbf{no person at all} on the project's rendered
skeletons, flat vector cel renders, or SD1.5 anime line-art output, while
reading a real photograph fine. YOLO11-pose does marginally better,
finding the cel figure in 2 of 6 tested clips at confidence 0.41 and 0.22
with $\sim$100~px mean keypoint error --- not ground truth, noise.
(Source: \texttt{ATTEMPTS.md} Attempt 19; corroborated by \secsym5.3's
later, broader finding that the same detectors read zero on Attempt 21's
material too, independent of cast size.) \textbf{One cause, two
symptoms}: the same off-manifold proportions that make ControlNet discard
the pose guide also make general human-pose estimators unable to read
this style at all.

\section{Experiments and Results II --- AnimateDiff: A Real Improvement That Still Fails the Gate}
\label{sec:results2}

Attempt 19 ended on a stated prediction: flicker is not a pose problem, so
fixing it needs an explicit temporal mechanism. \texttt{svg/animate.py} is
that mechanism --- AnimateDiff's motion module inserted into the same
frozen SD1.5 UNet, inference only, driven by the same COCO-18 guides.
Everything else from Attempt 19's final row is held fixed: \texttt{clip00}
(the wave), seed 1234, 20 steps, guidance 7.0, conditioning scale 1.0, 8
frames at 512~px. (Source for this entire section: \texttt{ATTEMPTS.md}
Attempt 20.) The guide's own inter-frame change at 512~px comes back as
1.09 --- identical to Attempt 19's final row --- so any change in the
ratio is entirely in the numerator.

\subsection{The first genuine reduction in flicker this project has measured, and it is not enough}
\label{sec:firstreduction}

\begin{table}[h]
\centering
\caption{AnimateDiff flicker ratios by resolution, 8 frames.
(Source: \texttt{ATTEMPTS.md} Attempt 20.)}
\label{tab:animatediffflicker}
\begin{tabular}{@{}L{6.5cm}rrr@{}}
\toprule
Method & Guide & Styled & Ratio \\
\midrule
A19 final, COCO-18, no temporal mechanism, 512~px & 1.09 & 39.58 & $\mathbf{36.2\times}$ \\
A20 + AnimateDiff motion module, 512~px, 8 frames & 1.09 & 27.32 & $\mathbf{25.0\times}$ \\
A20 + motion module, 384~px, 8 frames & 1.25 & 27.50 & $22.0\times$ \\
A20 + motion module, 256~px, 8 frames & 1.64 & 27.79 & $17.0\times$ \\
\bottomrule
\end{tabular}
\end{table}

This is the first intervention in the project's history to move flicker
\emph{down}: $36.2\times \rightarrow 25.0\times$, a 31\% cut in the
styled term, with the guide denominator unchanged. \textbf{Both halves of
this result matter and are reported together:} it is a real, measured
improvement, and --- the headline failure --- $\mathbf{25.0\times}$
\textbf{is still a hard fail} against the critic's $8\times$ bar;
\texttt{critic/evaluate.py} returns NOT-AN-ANIMATION at all three
resolutions.

\textbf{Trap: the resolution trend above is an artefact.} Styled flicker
is flat across resolution (27.32/27.50/27.79, a 1.7\% spread); the
apparent gain at lower resolution is the \emph{guide} moving more on a
smaller canvas ($1.09 \rightarrow 1.25 \rightarrow 1.64$), inflating the
ratio's denominator. \textbf{Only the 512~px row is comparable to Attempt
19; the $17.0\times$ figure must not be quoted as a result.}

\subsection{Warp error crossed a gate the flicker ratio hides}
\label{sec:warperror}

\begin{table}[h]
\centering
\caption{Flicker ratio vs.\ warp error. (Source: \texttt{ATTEMPTS.md} Attempt 20.)}
\label{tab:warpvsflicker}
\begin{tabular}{@{}L{4.3cm}L{3.5cm}L{4cm}l@{}}
\toprule
 & Flicker ratio & Warp error & Critic verdict \\
\midrule
A19 final, no motion module & $36.2\times$ (hard fail, bar $8\times$) & \textbf{0.185 --- hard fail}, bar 0.15 & NOT-AN-ANIMATION \\
A20 + motion module & $25.0\times$ (hard fail) & \textbf{0.135 --- soft fail}, pass bar 0.06 & NOT-AN-ANIMATION \\
\bottomrule
\end{tabular}
\end{table}

The motion module moved warp error out of the hard-fail band and into
the soft-fail band. The verdict does not change, because flicker alone
still disqualifies the clip.

\subsection{The motion module slightly hurts identity; an IP-Adapter anchor fixes it, and costs a little flicker back}
\label{sec:motionidentity}

\begin{table}[h]
\centering
\caption{Identity similarity across configurations. (Source: \texttt{ATTEMPTS.md} Attempt 20.)}
\label{tab:identitysim}
\begin{tabular}{@{}L{6.5cm}rr@{}}
\toprule
Run & Mean similarity & Worst similarity \\
\midrule
A19 final, no motion module & 0.934 & 0.787 \\
A20 + motion module & 0.897 & 0.744 \\
A20 + motion module + IP-Adapter anchor & \textbf{0.960} & \textbf{0.896} \\
\bottomrule
\end{tabular}
\end{table}

\texttt{svg/out/a20-ip/} adds \texttt{--ip-adapter --ip-scale 0.6},
anchored on the clip's own flat-colour vector render of frame 0.
Worst-case identity rises from 0.744 to 0.896, clearing the 0.45 pass
bar; mean rises to 0.960. \textbf{Caveat that must survive any reading of
this number:} the identity metric is an HSV colour-histogram embedding,
and the anchor run's own contact sheet shows its palette transferring
wholesale. What this result demonstrably locks is \textbf{colour scheme},
not identity in a structural sense. And the anchor does essentially
nothing for flicker ($25.0\% \rightarrow 25.3\%$, warp error $0.135
\rightarrow 0.142$, both marginally worse). \textbf{Two separate defects,
two separate mechanisms, and neither mechanism touches the other's
number} --- direct evidence for contribution 1 (\secsym1).

\subsection{Pose obedience got worse with the motion module, not better}
\label{sec:poseworse}

Attempt 19's corrected guide showed a raised wave arm clearly surviving
in 1 of 4 inspected frames. In the Attempt 20 sheet, the arm is raised in
\textbf{0 of 8}: the figure holds arms folded across its chest in every
frame. Temporal attention makes every frame agree with its neighbours ---
including agreeing on a pose that is not the guide's; joint denoising
averages the per-frame conditioning pull into a single clip-wide
compromise. \textbf{The motion module trades pose obedience for temporal
agreement} --- the second direct instance of contribution 1.

\subsection{Critic verdict: NOT-AN-ANIMATION throughout; one gate moved}
\label{sec:critic-a20}

Every configuration tried (8-frame headline run, both resolution
controls, the IP-Adapter run, the 16-frame run) returns
NOT-AN-ANIMATION, disqualified by flicker every time. The one gate that
moved is warp error, 0.185 (hard fail) $\rightarrow$ 0.135 (soft fail) at
8 frames.

\subsection{The stacked VRAM measurement the research survey could only mark unverified}
\label{sec:vram-stacked}

Resident weights, fp16, no offload: SD1.5 UNet + ControlNet-OpenPose +
AnimateDiff motion adapter = \textbf{4.41~GiB} of the 8~GiB card, before
a single latent exists.

\begin{table}[h]
\centering
\small
\caption{VRAM sweep, no offload, 20 steps. (Source: \texttt{ATTEMPTS.md} Attempt 20.)}
\label{tab:vramsweep}
\begin{tabular}{@{}lrrrrrr@{}}
\toprule
Config & $s$ & Peak alloc GiB & Peak reserved GiB & Guide & Styled & Ratio \\
\midrule
w512\_f16 & \textbf{OOM} & --- & 7.21 & --- & --- & --- \\
w448\_f16 & 173 & 6.46 & 7.04 & 1.23 & 99.02 & $80.5\times$ \\
w384\_f16 & 137 & 5.92 & 6.39 & 1.31 & 56.53 & $43.3\times$ \\
w320\_f16 & 101 & 5.46 & 5.78 & 1.54 & 59.71 & $38.9\times$ \\
w256\_f16 & 68 & 5.09 & 5.31 & 1.70 & 95.02 & $55.9\times$ \\
w512\_f8 & 129 & 5.74 & \textbf{6.15} & 1.09 & \textbf{27.32} & $\mathbf{25.0\times}$ \\
\bottomrule
\end{tabular}
\end{table}

\textbf{Only one configuration requires offload; everything else fits
resident.} \texttt{w448\_f16} peaks at 7.04~GiB reserved against
$\sim$7.5~GiB addressable. $512\text{px}\times16$ frames is the sole
exception (OOM at 7.21~GiB reserved); with
\texttt{enable\_model\_cpu\_offload()} it reaches 6.60~GiB / 207~s.
\textbf{Model-CPU-offload is very nearly free on this hardware} ---
$\sim$1.2~GiB saved at under 1.5--5\% wall-clock --- but this is reported
explicitly as hardware-specific: the 50\%-duty guard's mandatory idle
window absorbs the host-device copies, so the figure is not expected to
transfer to an ungoverned card.

\subsection{16 frames is worse than no motion module at all; the route closes here}
\label{sec:16frames}

At 512~px, flicker goes from $25.0\times$ at 8 frames to
$\mathbf{74.0\times}$ \textbf{at 16 frames}, with the guide denominator
essentially unchanged. Every 16-frame configuration lands between
$38.9\times$ and $80.5\times$, all worse than Attempt 19's $36.2\times$
with no motion module at all. The critic agrees more sharply: warp error
reaches 0.367 at 16 frames (back into the hard-fail band), worst-case
identity falls to 0.584, and \texttt{scene\_complexity} reports a
\textbf{modal cast of 0} --- no person detected in the plurality of
frames. \textbf{Below 384~px, the collapse is the base model's, not
AnimateDiff's}: \texttt{w256\_f16} contains no figure at all, only
coloured blocks, because SD1.5 was trained at 512~px and 256~px is out of
its distribution. This is the finding that closes the 16-frame route as
currently built: the VRAM work above establishes $512\times16$ is
technically reachable, and reaching it produces a clip roughly three
times less stable than the half-length one.

\subsection{Standing on the route}
\label{sec:standingroute-a20}

AnimateDiff is the first thing in this project's history to reduce
flicker rather than move it sideways or up, and it is not enough on its
own: a third off, still roughly $3\times$ over the bar, bought at the
price of the pose obedience Attempts 18 and 19 spent their entire budget
recovering, and it does not extend past 8 frames. \textbf{Flicker,
identity and pose obedience are three separate defects on this stack,
each needs its own mechanism, and the mechanisms interfere} --- the motion
module costs pose, the identity anchor costs a small amount of flicker
back. Nothing in the AnimateDiff route as built reaches USABLE, and the
project's own response to that closure is arc two (\secsym\secsym9--10).

\section{Experiments and Results III --- Multi-Character Scene Composition on CPU Alone (Attempt 21)}
\label{sec:results3}

Attempts 17--20 spent their whole budget trying to get styled anime out
of SD1.5 + ControlNet + AnimateDiff and closed that route by measurement
(\secsym\secsym7--8). Attempt 21 is a deliberate change of direction:
\textbf{several characters, a background and props, in one shot, driven
by a script} --- what the project's own log states has been the end goal
since Era 1, and had never been attempted. Style is explicitly not the
objective; legibility and composition are. The new module,
\texttt{scenescript.py}, uses no GPU at any point in generation --- the
whole pipeline is CPU-only numpy + Pillow + ffmpeg. (Source for this
entire section: \texttt{ATTEMPTS.md} Attempt 21.)

\subsection{The format}
\label{sec:sceneformat}

A scene script is a line-based text file declaring a cast, a background,
props, and a timeline of \texttt{<start> <end> <who> <action> [to X] [at
TARGET]} events. Actions are the rig's own
(walk/run/jump/kick/wave/dance/idle). Nothing new is drawn:
\texttt{stickman.StickFigure} supplies the joints, the same per-frame
equations Era 1 used render them, and \texttt{geovid.Scene} rasterises the
result.

\subsection{Five new problems, and what each turned out to be}
\label{sec:fiveproblems}

None of these exist when there is one figure doing one action.

\begin{enumerate}
\item \textbf{Root motion has to be taken away from the action.} Every
  action in \texttt{stickman.py} bakes its own absolute position; with
  three characters, every one teleports to the middle of the frame the
  moment it acts. The fix: the action supplies the pose, the script
  supplies the root --- each frame the compositor reads the action's hip,
  subtracts it, and re-anchors the joint set to where the script says the
  character is, keeping only the action's vertical hip offset so a walk's
  bob or a jump's arc survives while its horizontal sweep is discarded.
\item \textbf{Foot slide}, caused by (1): gait rate is derived from the
  distance the script actually asks for rather than assumed. Not solved
  exactly --- this matches average stride length, not per-foot contact ---
  and stated as approximate.
\item \textbf{Depth is not a free parameter, and occlusion is not a draw
  order.} A ground plane is defined once and depth defaults to scale.
  Because the cast is line art, draw order alone does not produce
  occlusion --- each character carries a matte, its own silhouette
  dilated by 4~px. The first version of that matte painted the backdrop
  colour directly into the frame, which erased scenery behind a character
  as well as the character behind it --- \textbf{a matte cannot be a
  colour.} The compositor was rebuilt to paste the original backdrop
  through the matte first (restoring scenery, deleting whatever was
  behind), then paste the character's colour through its ink mask.
\item \textbf{Props are stateful; characters are not.} A character is a
  pure function of $(\text{action}, \text{phase})$; a ball is not ---
  where it is at time $t$ depends on who kicked it and where it rolled.
  The ball is compiled into a segment list over the whole timeline (fly,
  then roll, then rest) and integrated, which is what lets it be kicked
  at one time and kicked back later from wherever it actually landed. The
  interaction runs the other way too: \texttt{kick ball at ana} reads the
  ball's actual current position and generates the kicker's approach to
  it --- the script never says where the kicker should stand.
\item \textbf{A background that is provably still.} The backdrop's
  layers are time-independent by construction and rendered exactly once,
  then copied as the base of every frame --- not ``stable to within a
  threshold,'' the same array.
\end{enumerate}

\subsection{Result 1 --- the videos}
\label{sec:a21result1}

\begin{table}[h]
\centering
\caption{Attempt 21 videos. (Source: \texttt{ATTEMPTS.md} Attempt 21.)}
\label{tab:a21videos}
\begin{tabular}{@{}L{4cm}L{3.5cm}rrrr@{}}
\toprule
File & Script & $s$ & Frames & Cast & Props \\
\midrule
\texttt{out/a21\_park.mp4} & \texttt{scenes/park\_meet.scene} & 12.5 & 375 & 3 & 7 \\
\texttt{out/a21\_street.mp4} & \texttt{scenes/street\_relay.scene} & 14.0 & 420 & 4 & 6 \\
\bottomrule
\end{tabular}
\end{table}

Render cost is $\sim$1.5~s of one CPU core per second of $854\times480$
video, and \textbf{no GPU is used to generate any of it} --- which after
Attempts 17--20 spent their entire budget inside a 7.5~GiB VRAM ceiling is
the point, not a footnote.

\subsection{Result 2 --- the critic}
\label{sec:a21result2}

\begin{table}[h]
\centering
\caption{Attempt 21 critic scores. (Identity values reflect the corrected
crop-regime figures from \secsym4.2/\secsym5.2; the original evaluation
reported \texttt{out/a21\_street.mp4} at identity worst 0.4368 and verdict
PROMISING, since corrected.)}
\label{tab:a21critic}
\begin{tabular}{@{}L{5cm}rrrrl@{}}
\toprule
Clip & warp\_err. ($\leq$0.06) & mot.\ cov.\ ($\geq$0.2\%) & id.\ worst & CLIPSIM & Verdict \\
\midrule
\texttt{out/a21\_park.mp4} & \textbf{0.0055} & \textbf{5.32\%} & \textbf{1.00} & 0.311 & \textbf{USABLE} \\
\texttt{out/a21\_street.mp4} & \textbf{0.0065} & 4.89\% & \textbf{1.00} (post-fix) & 0.336 & \textbf{USABLE} (post-fix) \\
\texttt{control\_wave.mp4} (ref.) & 0.0011 & 1.01\% & 1.00 & 0.329 & USABLE \\
\bottomrule
\end{tabular}
\end{table}

The composed multi-character scene holds the temporal gates the entire
diffusion route failed: warp error 0.0055 against the 0.06 pass bar and
the 0.15 hard-fail bar --- $\mathbf{25\times}$ \textbf{under the
hard-fail line} where Attempt 20's best AnimateDiff configuration sat at
0.135. Motion coverage is roughly $5\times$ the single-figure control's,
consistent with several characters moving independently. Flicker ratio is
N/A --- there is no conditioning guide, because nothing here is
conditioned on anything.

\subsection{Result 3 --- measuring the cast without a working detector}
\label{sec:a21result3}

\texttt{scene\_complexity} reports \texttt{mean\_count 0.00, modal\_count
0} on the park clip and \texttt{mean\_count 0.05, modal\_count 0} on the
street clip --- a photo-trained detector does not read flat line art at
any cast size, verified against this project's own single-figure positive
controls, which score the identical zero (\secsym4.6, \secsym5.3). So the
cast is measured another way, in \texttt{scenescript.py --audit}:
per-frame ink presence per declared character (colour match, connected
components).

\begin{table}[h]
\centering
\caption{Cast presence audit. (Source: \texttt{ATTEMPTS.md} Attempt 21.)}
\label{tab:a21audit}
\begin{tabular}{@{}lrl@{}}
\toprule
Clip & Characters & Frames with every character visible \\
\midrule
Park & 3 & \textbf{75/75 sampled} \\
Street & 4 & \textbf{84/84 sampled} \\
\bottomrule
\end{tabular}
\end{table}

Blob counts run 1.7--4.0 per character (a stick figure is normally two
components; the count rises exactly when another figure's matte cuts a
limb --- occlusion working, visible as a number).

\subsection{Result 4 --- the fresh-script test}
\label{sec:a21result4}

The log's standing criticism of this project is that its mechanisms only
ever work for the one example hand-authored with them. So
\texttt{scenes/street\_relay.scene} was written \textbf{after the engine
was frozen} and rendered with \textbf{no code change of any kind}: a
different cast size, background, props, action mix, and a three-way ball
relay instead of a two-way pass. It rendered first time and is coherent:
four distinguishable characters doing four distinct things in a shared
space, correct occlusion where they cross, the ball changing hands twice.
Critic (post-fix): warp 0.0065, motion 4.89\%, USABLE.

\subsection{What is honestly not solved}
\label{sec:a21notsolved}

Foot slide is fitted to average stride, not per-foot ground contact.
There is no author here, only an interpreter --- both scripts in this
attempt were written by hand, and Attempt 22 (\secsym10) is the response
to that gap specifically. No collision or contact: two characters can
walk through each other, and the ball's arc is kinematic. Actions cannot
blend across a clip boundary on one body --- a later-starting clip wins
outright. Staging is not checked: correct occlusion happily hides a
character behind a foreground prop for a whole clip, and two positions in
\texttt{park\_meet.scene} had to be moved by eye. CLIPSIM (0.311 park,
0.336 street) still cannot judge this content, for the same reason
established in \secsym4.3/\secsym5.1; \texttt{structural\_correctness},
the metric that \emph{did} separate correct from incorrect content on the
AniSVG path, is N/A here because a composed raster scene has no AniSVG
source --- a scene script is itself a declarative source and could in
principle be checked exactly the same way, but nothing does that yet.

\subsection{Standing}
\label{sec:a21standing}

The symbolic path now produces what this project has called the end goal
since Era 1 --- several characters, a background and props, in one
composed shot, from a script, as a real video, scoring USABLE from the
critic --- and it does it on the CPU. The two hard parts were neither the
rig nor the renderer: they were taking root motion away from the action,
and realising that a matte cannot be a colour. What is still missing is
not composition but authorship: no model writes these scripts, and
Attempt 21 by itself is subject to the same sentence as Era 1 --- the
best-looking output of the project is still hand-written. Attempt 22
(\secsym10) is the project's answer to exactly that gap.

\section{Experiments and Results IV --- Learning to Write Scene Scripts (Attempt 22)}
\label{sec:results4}

Attempt 21 removed half of the project's standing indictment --- the
composition machinery generalises to a script written after the engine
was frozen --- but left the other half standing: nothing writes the
scripts. This attempt trains a model to. (Source for this entire section:
\texttt{ATTEMPTS.md} Attempt 22.)

\subsection{The output space, stated before any model is trained}
\label{sec:outputspace}

The failure to avoid is Attempt 2's (Era 1): 100\% per-slot validation
accuracy on an output space of nine archetypes and four slots, a number
that ``meant almost nothing'' because the space was tiny.
\texttt{model/scene\_grammar.py} defines a genuinely combinatorial space
instead: cast size 1--5, character identity from 10 colours (distinct
within a scene), placement on a 21-point grid with continuous scale, 3
backgrounds, 0--3 props from 6 kinds, up to 12 timeline events with
continuous timing, 7 actions. A 3-character scene with 8 events already
has a cast assignment, 3 positions, 3 scales, a background, a prop list,
and 8 events each with two continuous times, an actor and an action ---
the sampler samples \emph{combinations}, not variants of a fixed script.

Two honest reductions stated up front: character names are canonical and
positional (\texttt{ana}, \texttt{bo}, \texttt{cy}, \texttt{dee},
\texttt{ell} by cast index --- a modelling simplification, not a claim
about the rest of the system), and numbers are emitted as integers in
tenths/hundredths rather than decimals --- a tokenizer economy that cut
mean target length from 118 tokens to 97.5.

\subsection{Training: t5-small, CPU, under nine minutes}
\label{sec:training-a22}

Same model and harness as Attempt 2 (\texttt{t5-small}, 60M params,
\texttt{model/train.py}), with three added CLI flags and no change to
its logic.

\begin{verbatim}
USE_TF=0 python model/train.py --train data/scene_train.jsonl \
    --val data/scene_val.jsonl --out scene_ckpt \
    --epochs 6 --batch 16 --lr 3e-4 --max-in 104 --max-out 176 --val-cap 200
\end{verbatim}

\textbf{8 CPU threads, 2250 steps, \texttt{train\_runtime} 530~s. No GPU
touched at any point.} Final eval\_loss 0.9240, against Attempt 2's
$\sim8\times10^{-5}$ --- \textbf{the expected result, not a regression}:
Attempt 2's target was a near-deterministic function of the prompt; here
the target --- every position, scale, timing, colour and decor choice ---
is mostly \emph{not} determined by the prompt at all, so cross-entropy
loss is not the metric that matters and is not quoted as a result. The
parse/render/match rates below are.

\subsection{In-distribution measurement: \texttt{model/scene\_eval.py}}
\label{sec:indist-a22}

150 held-out val prompts, disjoint from train, four rates: \textbf{parse}
(the string parses under strict \texttt{dsl\_to\_spec}), \textbf{valid}
(zero structural problems, before any repair), \textbf{render}
(\texttt{scenescript.py} builds and rasterises without raising, reported
after \texttt{repair()}), \textbf{prompt match} (per field, against only
the facts the prompt stated).

\begin{table}[h]
\centering
\caption{In-distribution rates, v1. (Source: \texttt{ATTEMPTS.md} Attempt 22.)}
\label{tab:v1indist}
\begin{tabular}{@{}lr@{}}
\toprule
 & \\
\midrule
Parse & \textbf{150/150 = 100.0\%} \\
Valid, zero problems, no repair & 103/150 = 68.7\% \\
Render through \texttt{scenescript.py} & \textbf{150/150 = 100.0\%} \\
Byte-identical to the reference DSL & \textbf{0/150 = 0.0\%} \\
Cast count & 145/145 = 100.0\% \\
\textbf{Every stated field correct at once} & \textbf{136/150 = 90.7\%} \\
\bottomrule
\end{tabular}
\end{table}

The 0.0\% exact-match line is worth reading twice: the model reproduces
the reference string never --- it picks its own positions, scales,
timings and scenery every time --- and still satisfies 90.7\% of the
prompts completely. The 31.3\% not structurally clean are dominated by
two faults, neither a parse failure: two characters sharing a colour (28
cases), overlapping kicks on one ball (10), an event ending after the
clip does (13). Both of the first two are \emph{global} constraints an
autoregressive, left-to-right decoder with no memory of what it has
already emitted has no mechanism to enforce; \texttt{repair()} fixes all
of them, which is why render is 100\% while valid is 68.7\%, and both
numbers are quoted rather than only the flattering one.

\subsection{The generalisation test --- the actual result}
\label{sec:generalisationtest}

\texttt{model/ood\_prompts.json} is 24 hand-written prompts, in six
groups, all outside the synthesis distribution: unseen phrasings and
verbs, unseen framings, unseen background wordings, cast sizes never seen
in training (six), action combinations never seen (a single character
kicking a ball --- the sampler only ever produced kicks with two or
more), five characters doing the same thing at once, and structurally
different prompts.

\begin{table}[h]
\centering
\caption{In-distribution vs.\ out-of-distribution, v1.
(Source: \texttt{ATTEMPTS.md} Attempt 22.)}
\label{tab:v1oodgap}
\begin{tabular}{@{}lrrr@{}}
\toprule
 & In dist. (150) & OOD (24) & Gap \\
\midrule
Parse & 100.0\% & \textbf{100.0\%} & 0.0 \\
Render & 100.0\% & \textbf{100.0\%} & 0.0 \\
Background & 100.0\% & \textbf{100.0\%} & 0.0 \\
\textbf{Cast count} & \textbf{100.0\%} & \textbf{33.3\%} & \textbf{66.7} \\
\textbf{Every stated field at once} & \textbf{90.7\%} & \textbf{25.0\%} & \textbf{65.7} \\
\bottomrule
\end{tabular}
\end{table}

\textbf{The headline number of Attempt 22 is a 65.7-point compositional
generalisation gap, and essentially all of it is one failure.} The
\emph{format} generalises perfectly --- 24/24 parse, 24/24 render, 24/24
correct background. What does not generalise is counting the cast: 16 of
24 failures are \texttt{cast\_count}, and 15 of those emit a cast of
exactly \textbf{one}.

\subsection{The cause, isolated by a controlled probe}
\label{sec:causeprobe}

The obvious reading --- ``it cannot count past what it saw'' --- is
wrong, and a controlled probe says so. \texttt{model/eval\_out/clause\_probe.json}
holds the stated numeral and the number of described characters
independent and asks for the cast size:

\begin{table}[h]
\centering
\caption{Clause probe, v1: emitted cast size (bold = numeral and clause
count agree). (Source: \texttt{ATTEMPTS.md} Attempt 22.)}
\label{tab:clauseprobe1}
\begin{tabular}{@{}lrrrrr@{}}
\toprule
Stated numeral & 1 clause & 2 clauses & 3 clauses & 4 clauses & 5 clauses \\
\midrule
``one people\ldots'' & \textbf{1} & & & & \\
``two people\ldots'' & 1 & \textbf{2} & & & \\
``three people\ldots'' & 1 & 2 & \textbf{3} & & \\
``four people\ldots'' & 1 & 2 & 3 & \textbf{4} & \\
``five people\ldots'' & 1 & 2 & 3 & 4 & \textbf{5} \\
\bottomrule
\end{tabular}
\end{table}

\textbf{Emitted cast size equals the number of described characters in
15 of 15 probes, and the stated numeral is ignored in all 10
disagreements.} The model can produce a cast of five --- the capability
is there. It never learned to read the numeral; it learned to count
clauses.

\textbf{And the cause of \emph{that} is in the data, not the model.}
Measured on the training set: the number of characters the prompt
describes equals cast size in \textbf{5681 of 5688 = 99.9\%} of
synthesized examples, because the synthesiser gave every character its
own clause. ``Count the clauses'' and ``read the numeral'' were the same
function on 99.9\% of the training distribution, and gradient descent
took the cheaper one. This is a spurious-correlation failure the project
built into its own dataset, and every out-of-distribution prompt that
states a count without describing each character individually --- which
is how people actually write --- walks straight into it.

\subsection{Removing the confound and retraining}
\label{sec:deconfound}

\texttt{scene\_synth.py} gained a \texttt{--deconfound} flag: with $p=0.45$
a multi-member cast's prompt states the count but describes only a random
subset of the characters; with $p=0.22$ the whole cast is described in
one clause sharing one action. Nothing else changed --- same grammar,
model, hyperparameters, seed. Measured on the new training set, clause
count equals cast size in \textbf{68.0\%} of examples, down from 99.9\%.
Retrained: 6 epochs, \texttt{train\_runtime} 531.7~s, same 8 CPU threads,
final eval\_loss 0.9470 (v1: 0.9240 --- slightly \emph{worse} loss on a
harder, less predictable dataset, again the point about loss not being
the result).

The same probe on v2: \textbf{14 of 15 agree with the numeral, against 5
of 15 for v1} (the one miss is an ungrammatical probe string). The
diagnosis was correct and it was a data fault, not a model fault.

\begin{table}[h]
\centering
\caption{v1 vs.\ v2, in-distribution and out-of-distribution.
(Source: \texttt{ATTEMPTS.md} Attempt 22.)}
\label{tab:v1v2gap}
\begin{tabular}{@{}lrrrr@{}}
\toprule
 & v1 in-dist & v1 OOD & v2 in-dist & v2 OOD \\
\midrule
Parse & 100.0\% & 100.0\% & 100.0\% & 100.0\% \\
Valid, no repair & 68.7\% & 62.5\% & 40.0\% & 54.2\% \\
\textbf{Cast count} & \textbf{100.0\%} & \textbf{33.3\%} & \textbf{100.0\%} & \textbf{83.3\%} \\
Background & 100.0\% & 100.0\% & 100.0\% & 100.0\% \\
\textbf{All stated fields} & \textbf{90.7\%} & \textbf{25.0\%} & \textbf{86.0\%} & \textbf{62.5\%} \\
\textbf{Generalisation gap} & --- & \textbf{65.7 pts} & --- & \textbf{23.5 pts} \\
\bottomrule
\end{tabular}
\end{table}

The gap is cut from 65.7 points to 23.5 points, at a measured cost of 4
points of in-distribution accuracy ($90.7\% \rightarrow 86.0\%$) ---
\textbf{reported as a cost, not hidden.} Two things survive deconfounding
and are stated plainly:

\begin{itemize}
\item \textbf{Six characters is a hard capability limit, not a phrasing
  problem.} All three \texttt{cast\_size} prompts fail on both models; v2
  answers ``six stick figures'' with a cast of 4, ``half a dozen'' with
  3, ``a crowd of six'' with 4. The training distribution contains 1--5
  and the model does not extrapolate to 6 --- it saturates.
\item \textbf{The \texttt{valid} rate got \emph{worse}, $68.7\%
  \rightarrow 40.0\%$, because of duplicate colours} (71 of 150 v2
  outputs against 28 of 150 for v1). Group scenes push mean cast size up,
  and ``no two characters share a colour'' is a global constraint an
  autoregressive decoder cannot see. Rendering is unaffected (repair
  recolours), but this is a real defect \texttt{repair()} is currently
  carrying, not one that disappeared.
\item \textbf{Positions still collapse to a handful of modal values} (the
  deliverable below places two characters at 9.6 and 9.2 on a
  ten-unit-wide frame). The prompt says nothing about where anyone
  stands, so this slot has no supervision at all, and the decoder answers
  an unasked question with its prior.
\end{itemize}

\subsection{Does it survive the critic? A rate, not an anecdote}
\label{sec:survivecritic}

Ten generated scenes (five in-distribution, five out-of-distribution)
rendered through the unmodified \texttt{scenescript.py} and scored:

\begin{table}[h]
\centering
\caption{Ten generated scenes, critic scores. (Source: \texttt{ATTEMPTS.md} Attempt 22.)}
\label{tab:tengenerated}
\begin{tabular}{@{}L{5cm}rrrl@{}}
\toprule
Clip & Warp ($\leq$0.06) & Mot.\ cov.\ ($\geq$0.2\%) & Id.\ worst & Verdict \\
\midrule
a22\_batch\_ind0--ind4 (5 clips) & 0.0013--0.0042 & 1.45\%--3.25\% & 1.00 & USABLE (all 5) \\
a22\_batch\_ood0, ood2, ood3, ood4 & 0.0005--0.0029 & 0.78\%--2.51\% & 1.00 & USABLE (all 4) \\
a22\_batch\_ood1 & 0.0002 & \textbf{0.15\%} & 1.00 & \textbf{NOT-AN-ANIMATION} \\
\bottomrule
\end{tabular}
\end{table}

\textbf{Critic pass rate: 9/10 USABLE.} The one failure is worth more
than the nine passes. The prompt was \emph{``an empty studio, two
figures, no scenery: they simply stand around doing nothing''} --- one of
the project's own out-of-distribution prompts --- and the model answered
it exactly right: two characters, studio background, no props, four
\texttt{idle} events and nothing else. Motion coverage 0.15\% against the
critic's 0.2\% floor fires the gate. \textbf{The model obeyed the prompt
and the judge rejected the obedience} --- the motion floor doing
precisely what it was added to do (\secsym4.1) meeting, for the first
time in this project, a clip that is \emph{supposed} to be nearly still.
This is not recorded as a model error. \texttt{identity\_preservation}
reads 1.00 on every clip because the crop-regime fix (\secsym5.2) now
forces the fallback crop consistently when YOLO cannot detect the
material, and \texttt{scene\_complexity} reports \texttt{UNMEASURABLE}
rather than a false cast of zero (\secsym4.6) --- both are Attempt 21
findings the evaluation harness has since acted on.

\subsection{The deliverable: an unseen prompt to an .mp4}
\label{sec:deliverable-a22}

\begin{verbatim}
python model/scene_infer.py \
  "two friends meet in the park, one waves, then they kick a ball around" \
  --ckpt model/scene_ckpt2 -o scenes/a22_generated.scene
python scenescript.py scenes/a22_generated.scene \
  -o out/a22_generated.mp4 --frames out/a22_generated_frames
\end{verbatim}

The model emitted, with \textbf{zero structural problems and zero
repairs}, a scene that rendered to \texttt{out/a22\_generated.mp4}:
9.1~s, 273 frames, 2 characters, 2 props, 4 events. Critic: warp error
\textbf{0.0022}, motion coverage 2.03\%, identity 1.00, CLIPSIM 0.297,
\textbf{USABLE}. Cast audit: both characters visible in 55/55 sampled
frames. A second deliverable, from an out-of-distribution prompt with a
three-way group/pair structure (``three friends in the park: all three
dance at the same time, then two of them pass a ball while the third one
just watches''), also rendered: 9.2~s, 3 characters, 7 events, warp
0.0046, motion 3.74\%, \textbf{USABLE}. \textbf{These are the first
videos in this project whose script was not written by a human.}

\section{Experiments and Results V --- Staging Is Mostly a Decoding Artefact (Attempt 23)}
\label{sec:results5}

Attempt 22 ended by naming staging as ``the single biggest defect in
every generated scene'' and the thing standing between the pipeline and
a scene that looks composed. Attempt 23 takes that sentence apart, and
does so against \texttt{paper/RESEARCH2.md} (\secsym2), whose \secsym7
asked for a falsifiable prediction per technique stated \emph{before} the
measurement --- those predictions are scored honestly below, including
the ones that were wrong. (Source for this entire section:
\texttt{ATTEMPTS.md} Attempt 23.)

\subsection{Versioning and reproducibility, first}
\label{sec:versioning-a23}

Before any new code ran, v1 and v2 were copied to
\texttt{model/checkpoints/\{v1-confounded,v2-deconfounded\}/} alongside
the exact data that trained them, and \texttt{model/verify\_checkpoints.py}
confirmed both still load and reproduce their documented behavior ---
including v2-deconfounded's byte-for-byte reproduction of Attempt 22's
published output (\secsym6). v3 trains to a new directory; nothing
already preserved is touched.

\subsection{A staging metric, independently calibrated}
\label{sec:stagingmetric}

\texttt{model/scene\_staging.py} samples 61 instants across a clip and,
at each, reads every character's actual position and the half-width of
its \emph{current} silhouette --- built off the rig by posing
\texttt{stickman.StickFigure} through all seven actions (an idle figure
is 0.34 units wide about its hip, a kicking one 0.92, nearly three times
as much), not assumed as a single constant. Overlap only counts as a
staging fault when two silhouettes overlap \emph{and} the characters
stand on ground planes within 0.20 units of each other --- a near
character crossing a far one is depth, not a defect, and is exactly what
Attempt 21's occlusion matte was built to render correctly. The composite
\texttt{staging} score is the mean of three sub-scores (pairwise
separation adequacy, stage-span fraction, occlusion fraction) and is
undefined (\texttt{NaN}) for a one-character scene, since a solo cast has
no staging to measure.

\textbf{Two candidate sub-terms were built, measured against the
project's own reference clips, and thrown out because they disagreed
with clips already known to look good} --- reported here because the
decision itself is part of the metric's validation: an edge-clearance
term scored 0.00 on both hand-written scenes (one character's hand
crosses the frame edge for a few frames while walking to a
script-specified position) and was demoted to a diagnostic; a
\texttt{min()}-over-the-clip separation term scored the hand-written
\texttt{street\_relay.scene} at 0.016 because two characters pass within
a fraction of a second during a scripted ball hand-off, and was replaced
by a mean-adequacy version. \textbf{The honest reading: a metric that
disagrees with the clips it is calibrated against is wrong, and the fix
is to change the metric, not re-describe the clips.}

\textbf{Calibration}, against the project's two hand-written reference
scenes (the best-looking output of the whole project, per Attempts
21--22) and eleven of Attempt 22's model-written scenes:

\begin{table}[h]
\centering
\caption{Staging metric calibration. (Source: \texttt{ATTEMPTS.md} Attempt 23 \secsym2.)}
\label{tab:stagingcalib}
\begin{tabular}{@{}L{7cm}r@{}}
\toprule
 & Staging \\
\midrule
\texttt{scenes/park\_meet.scene} (hand-written, 3 cast) & \textbf{1.000} \\
\texttt{scenes/street\_relay.scene} (hand-written, 4 cast) & \textbf{0.825} \\
\textbf{Hand-written mean} & \textbf{0.912} \\
Best model-written (\texttt{a22\_batch\_ood1}) & 0.747 \\
\texttt{scenes/a22\_generated.scene} (the Attempt 22 deliverable) & 0.332 \\
Worst model-written (\texttt{a22\_batch\_ood4}) & 0.000 \\
\textbf{Model-written mean, $n=11$} & \textbf{0.352} \\
\bottomrule
\end{tabular}
\end{table}

\textbf{The separation is complete: the worst hand-written scene (0.825)
scores above the best model-written one (0.747), with no overlap.} The
Attempt 22 deliverable's 0.332 decomposes as span 0.346 (two characters
spanning 15.6\% of the usable stage) and occlusion 0.000 (they cover each
other in 34.4\% of sampled observations) --- the 9.6/9.2 stack, as a
number, for the first time.

\subsection{The root cause is not what was expected: the training data is fine}
\label{sec:rootcause}

The obvious next move --- the synthesizer must be emitting stacked scenes
--- was checked before anything was changed. Scoring 400 rows of the
actual v2 training data with the new metric:

\begin{table}[h]
\centering
\caption{Training-data staging. (Source: \texttt{ATTEMPTS.md} Attempt 23 \secsym3.)}
\label{tab:trainingstaging}
\begin{tabular}{@{}L{7cm}r@{}}
\toprule
 & Staging \\
\midrule
Hand-written scenes & 0.912 \\
\textbf{v2 training targets ($n=371$ multi-cast of 400)} & \textbf{0.865 mean, 0.908 median} \\
v2 model outputs & 0.352 \\
\bottomrule
\end{tabular}
\end{table}

\textbf{The training data is not badly staged} --- it sits 0.05 below the
hand-written clips and 0.51 above what the model trained on it produces
--- so the hypothesis this attempt was scheduled to act on is false. The
actual mechanism, named but not quantified in Attempt 22 \secsym10, is
\textbf{conditional mode collapse}: the prompt almost never states where
anyone stands, position slots are sampled independently in the original
synthesizer, so the training distribution's conditional structure between
position slots is close to uniform, and a near-uniform conditional has no
informative argmax --- beam search then lands on whatever small marginal
bias exists, for every slot at once. This reframes the fix: not ``produce
nicer training scenes'' but ``make position slots statistically depend on
each other,'' which is a decoding and data-conditioning question, not a
data-quality one.

\subsection{What was built, and the predictions made before measuring}
\label{sec:whatbuilt}

Per RESEARCH2 \secsym7's own requirement, each change carries a stated
prediction. \texttt{model/scene\_ckpt2/config.json} settles one item
immediately, without a training run: \texttt{d\_model} is \textbf{512}
against closed vocabularies of at most ten words (10 colours, 8 names, 3
backgrounds, 7 actions), which ``When Can Transformers Count to $n$?''
\cite{yehudai2024counttransformers} predicts is far above the regime
where exact counting becomes numerically unstable --- ruling out ``make
the model wider'' as a fix before it is tried, and confirmed against the
actual checkpoint rather than t5-small's published defaults. The
remaining changes: a \texttt{UniqueColourProcessor} logits processor
masking already-used colour tokens (inference-only, $\sim$60 lines,
hand-written against \texttt{transformers.LogitsProcessor} rather than
adopting Outlines/Guidance per RESEARCH2 \secsym4's own recommendation
for one small constraint); a \texttt{space\_out()} deterministic
post-decode spacing pass enforcing a 1.25-unit minimum separation between
emitted $x$-coordinates while preserving the model's chosen left-right
order; positions sampled \emph{jointly} rather than independently in the
synthesizer, with 12\% of multi-cast scenes deliberately clustered and
described as such; spatial-language paraphrases (seven fact kinds ---
leftmost, rightmost, middle, back, front, spread, close) attached to a
character's own action clause, used only when true of the sampled scene;
and cast size extended from 1--5 to 1--8. \textbf{RESEARCH2's own
top-ranked staging fix --- the deterministic spacing pass --- and its own
\secsym2 diagnosis of ``conditional mode-collapse'' were both followed;
the survey did not consider the decoder itself as a candidate cause},
which \secsym11.6 below shows mattered.

\textbf{The DETR-style Hungarian set-decoder (RESEARCH2 \secsym4), the
one surveyed technique that structurally fixes staging and duplicate
colours together, was deliberately not built.} It is the only item
requiring genuine architecture work (a new decoder head, a new
order-invariant matching loss) against a nine-minute CPU retraining
budget, and --- as \secsym11.8 and \secsym11.12 below show --- a
decoding-argument change and a rerank pass closed most of the gap without
it. It remains the principled fix and is future work (\secsym13), not an
unexplained omission.

One more measurement caught a defect that would have quietly corrupted
this attempt: eight-character casts need up to eighteen timeline events,
and at Attempt 22's \texttt{--max-out 176}, 11.28\% of the new training
targets would have been silently truncated mid-scene. v3 therefore
trains at \texttt{--max-out 256} (0\% truncated, measured over all 6000
rows) and \texttt{--max-in 160} --- a hyperparameter difference from v2
forced by the data, not chosen, and noted wherever the two are compared.

\subsection{v3 training: loss rose again, for the third time and the same reason}
\label{sec:v3training}

\begin{verbatim}
python model/train.py --train model/data/scene3_train.jsonl \
  --val model/data/scene3_val.jsonl --out model/scene_ckpt3 \
  --epochs 6 --batch 16 --lr 3e-4 --max-in 160 --max-out 256 --val-cap 200
\end{verbatim}

\textbf{8 CPU threads, 2250 steps, \texttt{train\_runtime} 548.5~s. No
GPU at any point except the critic's YOLO pass.}

\begin{table}[h]
\centering
\caption{v1/v2/v3 loss and runtime. (Source: \texttt{ATTEMPTS.md} Attempt 23 \secsym5.)}
\label{tab:v3loss}
\begin{tabular}{@{}lrrr@{}}
\toprule
 & v1 & v2 & v3 \\
\midrule
Final eval\_loss & 0.9240 & 0.9470 & \textbf{1.0042} \\
train\_runtime & 530~s & 531.7~s & \textbf{548.5~s} \\
\bottomrule
\end{tabular}
\end{table}

Loss has risen at every step of this sequence, and every rise bought a
capability: v3's target space has eight possible cast sizes instead of
five and up to eighteen events instead of twelve, raising the entropy of
``the correct answer'' without making the model worse. Loss is not the
result and is not quoted as one, for the third time in this project.

\subsection{The result that reframes the attempt}
\label{sec:reframes}

Reported in full in \secsym5.4: decoding the same checkpoint on the same
prompts under beam search versus sampling moves staging from 0.188 to
0.862 on v3 and 0.333 to 0.845 on v2, with distinct emitted $x$-values
going from 3 to 55 on v3. \textbf{The position collapse Attempt 22
attributed to the model was mostly the decoder.}

\subsection{The four-cell ablation: model, decoder, or a pass that fixes it afterwards}
\label{sec:fourcell}

All cells: the same 150 in-distribution and 24 out-of-distribution
prompts, \texttt{beam4}. \texttt{+uc} = the colour processor.

\begin{table}[h]
\centering
\footnotesize
\caption{Four-cell ablation, in distribution (150 prompts).
(Source: \texttt{ATTEMPTS.md} Attempt 23 \secsym7;
\texttt{model/eval\_out3/compare\_indist.txt}.)}
\label{tab:fourcellindist}
\begin{tabular}{@{}L{5cm}rrrr@{}}
\toprule
 & v2 & v2+uc & v3 & v3+uc \\
\midrule
Parse & 100.0 & 100.0 & 98.7 & 98.0 \\
\textbf{Valid, no repair} & 40.0 & \textbf{72.0} & 56.0 & \textbf{80.7} \\
\textbf{Duplicate-colour problems} & \textbf{71} & \textbf{0} & \textbf{46} & \textbf{0} \\
Actions per character & 88.7 & 90.0 & 85.8 & 85.7 \\
Layout (stated spatial fact obeyed) & n/a & n/a & 44.3 & 45.8 \\
\textbf{ALL, layout excluded (like-for-like)} & \textbf{86.0} & \textbf{88.0} & \textbf{85.8} & \textbf{85.7} \\
\textbf{Staging} & 0.360 & 0.363 & \textbf{0.178} & 0.182 \\
\textbf{Staging, + spacing pass} & \textbf{0.670} & 0.670 & 0.587 & 0.585 \\
Declared min.\ separation & 0.157 & 0.180 & 0.018 & 0.018 \\
Declared min.\ separation, + pass & \textbf{1.217} & 1.218 & 1.227 & 1.228 \\
Min.\ separation, whole clip & 0.069 & 0.074 & 0.010 & 0.010 \\
Min.\ separation, whole clip, + pass & 0.348 & 0.343 & 0.238 & 0.245 \\
Staging $\geq0.825$, + pass & 19.3\% & 17.9\% & 17.0\% & 15.7\% \\
\bottomrule
\end{tabular}
\end{table}

\textbf{The honest decomposition of staging, in distribution:} what the
model learned, read by beam search --- 0.178; plus a deterministic pass
enforcing it afterward --- 0.587; what the model learned, read by
sampling instead --- \textbf{0.862}; hand-written reference clips ---
0.825--1.000. \textbf{Under beam search the deterministic pass does
essentially all of the work, and that is a good engineering result and a
bad science claim} --- it raises declared minimum separation from 0.018
to 1.227 by construction (it cannot fail to) and moves staging to 0.587,
but never reaches the hand-written band, and only 17\% of scenes clear
the worse reference clip. Changing the decoder alone reaches 0.862, past
that clip, because it is the model's own answer rather than an imposed
one --- \textbf{the pass is a floor, not a fix.}

\textbf{Two smaller findings inside the ablation, both falsifying part of
RESEARCH2's own stated prediction 2.} The spacing pass fixes
\emph{declared} minimum separation exactly as designed ($0.157
\rightarrow 1.217$) but \emph{over-clip} minimum separation only reaches
$0.069 \rightarrow 0.348$, because the pass rewrites the cast section and
never touches \texttt{walk}/\texttt{run} destinations in the timeline ---
a character correctly spaced at $t{=}0$ can still walk into somebody.
RESEARCH2 \secsym7 prediction 2 said minimum pairwise distance rises to
an enforced floor with no further qualification; \textbf{as stated, that
is falsified} --- true of the declared positions, false of the clip,
because nothing had measured over-clip separation before this attempt
built the metric to do so. Separately, v3's 98.7\% parse rate (in
distribution) is a length cap, not a format failure: both parse failures
end mid-token at seven- and eight-character casts whose timelines exceed
\texttt{--max-out 256} at inference --- the same class of bug as the
training truncation caught in \secsym11.4, one layer further out.

\subsection{Constrained decoding: the cleanest result in the attempt}
\label{sec:constraineddecoding}

\begin{table}[h]
\centering
\small
\caption{Constrained decoding results. (Source: \texttt{ATTEMPTS.md} Attempt 23 \secsym8.)}
\label{tab:constraineddecoding}
\begin{tabular}{@{}L{5.2cm}rrrrr@{}}
\toprule
 & v1 (A22) & v2 & v2 + proc. & v3 & v3 + proc. \\
\midrule
In-dist dup.-colour problems & 28/150 & 71/150 & \textbf{0/150} & 46/150 & \textbf{0/150} \\
In-dist valid, no repair & 68.7\% & 40.0\% & \textbf{72.0\%} & 56.0\% & \textbf{80.7\%} \\
OOD dup.-colour problems & --- & 7/24 & \textbf{0/24} & 6/24 & \textbf{0/24} \\
OOD valid, no repair & 62.5\% & 54.2\% & \textbf{58.3\%} & 50.0\% & \textbf{70.8\%} \\
\bottomrule
\end{tabular}
\end{table}

\textbf{The constraint is enforced exactly: 0 duplicate colours in 348
generated scenes across four cells, by construction rather than by
rate.} The valid-without-repair rate Attempt 22 reported as its worst
regression ($68.7\% \rightarrow 40.0\%$, v1$\rightarrow$v2) is not only
recovered but beaten: \textbf{80.7\%} on v3 with the processor, the
highest raw validity rate this project has recorded at that point.
RESEARCH2's own prediction --- duplicates go to 0 with zero change to any
other category --- is \textbf{half right}: the count goes to exactly
0/150 and 0/24 as predicted, but masking a token shifts the beam path,
moving v2's \texttt{ALL} up ($86.0 \rightarrow 88.0$) and v3's parse rate
down ($98.7 \rightarrow 98.0$, one extra length-cap truncation).
``Orthogonal'' was the right intuition and the wrong word: the constraint
is exact, its side effects are real but second-order.

\subsection{What sampling costs, on all 174 prompts}
\label{sec:samplingcosts}

Both models were re-evaluated end to end with \texttt{do\_sample=True,
top\_p=0.95, temperature=1.0}, colour processor on, same 150 + 24
prompts as every other cell.

\begin{table}[h]
\centering
\caption{Sampling costs, in distribution, v3.
(Source: \texttt{ATTEMPTS.md} Attempt 23 \secsym9.)}
\label{tab:samplingcosts}
\begin{tabular}{@{}L{5.5cm}rrr@{}}
\toprule
In distribution, v3 & beam4 & sampled & Delta \\
\midrule
Valid, no repair & \textbf{80.7\%} & 60.7\% & $\mathbf{-20.0}$ \\
Actions per character & \textbf{85.7\%} & 72.5\% & $\mathbf{-13.2}$ \\
Layout & \textbf{45.8\%} & 25.3\% & $\mathbf{-20.5}$ \\
\textbf{Staging} & 0.182 & \textbf{0.875} & $\mathbf{+0.693}$ \\
\textbf{Scenes clearing 0.825} & \textbf{0.0\%} & \textbf{74.3\%} & $\mathbf{+74.3}$ \\
\bottomrule
\end{tabular}
\end{table}

\textbf{In distribution, sampling buys 0.693 of staging for 13--20 points
of obedience}, concentrated exactly in the fields the prompt \emph{does}
specify (\texttt{actions\_per\_character}, \texttt{layout});
\texttt{cast\_count}, read from a numeral rather than a spatial
judgement, does not move (100.0\% either way). \textbf{Out of
distribution it is not a trade at all} --- \texttt{ALL} stated fields
rises $70.8\% \rightarrow 75.0\%$ at the same time staging rises $0.099
\rightarrow 0.804$, because on unfamiliar input the beam-search mode was
never a better answer, only a more confident one. \textbf{\texttt{valid}
is where sampling genuinely costs something real}: $80.7\% \rightarrow
60.7\%$ in distribution, because unconstrained sampling occasionally
emits a token outside the DSL entirely (\texttt{yellow}, not one of the
ten trained colours) that the strict parser correctly rejects --- the
strongest argument in this attempt for extending the colour processor's
approach to the whole grammar, per RESEARCH2 \secsym4's own fallback
recommendation once sampling replaces beam search.

\subsection{Cast extrapolation: the wall moved from six to nine, exactly as predicted}
\label{sec:castextrapolation}

\texttt{model/scene\_probe.py}, the Attempt 22 clause probe extended to
numeral 9:

\begin{table}[h]
\centering
\caption{Cast extrapolation, clause probe extended to 9.
(Source: \texttt{ATTEMPTS.md} Attempt 23 \secsym10.)}
\label{tab:castextrapolation}
\begin{tabular}{@{}lrr@{}}
\toprule
Range & v2 (trained 1--5) & v3 (trained 1--8) \\
\midrule
1--5, the Attempt 22 range & 14/15 & \textbf{14/15} \\
6--8, new in the v3 range & \textbf{0/15} & \textbf{15/15} \\
9, one step past the trained ceiling & 0/5 & \textbf{0/5} \\
\bottomrule
\end{tabular}
\end{table}

v3 answers ``nine people'' with a cast of \textbf{6, every single time,
at every clause count} --- a clean ceiling, not random failure, exactly
the pattern v2 showed at 4--5 when asked for six. On the 24 Attempt 22
out-of-distribution prompts, the \texttt{cast\_size} group goes 0/3 (v1)
$\rightarrow$ 0/3 (v2) $\rightarrow$ \textbf{3/3 (v3)}, and OOD
\texttt{cast\_count} goes $33.3\% \rightarrow 83.3\% \rightarrow
\mathbf{95.8\%}$. \textbf{RESEARCH2's prediction 3 is confirmed in both
halves} --- sizes inside the new range reach full parity, and the size
one step past the new ceiling reproduces the same saturation pattern
rather than failing randomly --- and per RESEARCH2 \secsym3's own reading
of the compositional-generalization literature (SCAN/GECA: what fixes
systematic generalization is recombination inside or near the trained
distribution, not extrapolation past it), \textbf{we report this as
confirming a known general seq2seq limitation, not as a defect specific
to this project.} Training on 1--12 would, by the same pattern, move the
wall to 13.

\subsection{The spatial supervision did not work, and the control is what shows it}
\label{sec:spatialfail}

v3 obeys the spatial fact its own prompt states in 44.3\% of
in-distribution scenes --- a number that reads like a capability until
it is compared against a control. \texttt{model/scene\_layout\_baseline.py}
scores each prompt's spatial facts against a \textbf{different} prompt's
emitted scene, averaged over every mismatched pair: with a cast of two,
``leftmost'' is true about half the time whatever the model does.

\begin{table}[h]
\centering
\footnotesize
\caption{Spatial-language lift over a mismatched-pair control.
(Source: \texttt{ATTEMPTS.md} Attempt 23 \secsym11.)}
\label{tab:spatiallift}
\begin{tabular}{@{}L{6.3cm}rrrr@{}}
\toprule
Cell & $n$ & Layout obeyed & Mismatch chance & Lift \\
\midrule
v3 beam4, in distribution & 97 & 44.3\% & 41.7\% & $\mathbf{+2.7}$ \\
v3 sampled, in distribution & 99 & 25.3\% & 27.3\% & $\mathbf{-2.0}$ \\
v3 beam4, hand-written spatial OOD & 12 & 50.0\% & 49.2\% & $+0.8$ \\
v3 sampled, hand-written spatial OOD & 12 & 58.3\% & 34.8\% & $+23.5$ \\
v2 beam4, hand-written spatial OOD (control) & 12 & 58.3\% & 45.5\% & $\mathbf{+12.9}$ \\
\bottomrule
\end{tabular}
\end{table}

\textbf{In distribution the lift is $+2.7$ points and, under sampling,
$-2.0$: the model is not obeying the spatial instruction, it is producing
geometry that happens to satisfy it at the base rate}, despite 61.2\% of
v3's training prompts stating a true spatial fact. The one apparently
positive cell (v3 sampled on hand-written spatial prompts, $+23.5$) does
not survive its own control: v2, which never saw a spatial phrase in
training, scores $+12.9$ on the identical prompts. With $n=12$ and a
never-trained model reaching half the lift, that cell is noise.
\textbf{This falsifies half of the attempt's own prediction 3}
(RESEARCH2 \secsym7): the joint position sampling half demonstrably
worked (\secsym11.6, \secsym11.9), the spatial-language half bought no
measurable obedience --- likely because binding an ordinal phrase to a
specific cast member is a harder problem than reading a numeral, and
3,673 training examples on a 60M model trained for nine minutes was not
enough of it. \textbf{Every staging result in this attempt comes from
the decoder and the position sampler, not from the prompt following
spatial instructions.}

\subsection{Sample-and-rerank: the configuration that keeps both}
\label{sec:samplerank}

\texttt{SceneWriter(rerank=4)} draws four samples and keeps the first
with the fewest structural problems, tie-breaking on draw order (not
staging score) so the staging number stays an unbiased read rather than a
selected one.

\begin{table}[h]
\centering
\caption{Sample-and-rerank, in distribution.
(Source: \texttt{ATTEMPTS.md} Attempt 23 \secsym12.)}
\label{tab:samplerankindist}
\begin{tabular}{@{}L{5cm}rrr@{}}
\toprule
In distribution & beam4 & sampled & sampled + rerank(4) \\
\midrule
\textbf{Valid, no repair} & 80.7\% & 60.7\% & \textbf{94.0\%} \\
Actions per character & \textbf{85.7\%} & 72.5\% & 72.7\% \\
\textbf{Staging} & 0.182 & 0.875 & \textbf{0.879} \\
\textbf{Scenes $\geq0.825$} & 0.0\% & 74.3\% & \textbf{75.2\%} \\
\bottomrule
\end{tabular}
\end{table}

\begin{table}[h]
\centering
\caption{Sample-and-rerank, out of distribution.
(Source: \texttt{ATTEMPTS.md} Attempt 23 \secsym12.)}
\label{tab:samplerankood}
\begin{tabular}{@{}L{5cm}rrr@{}}
\toprule
Out of distribution & beam4 & sampled & sampled + rerank(4) \\
\midrule
\textbf{Valid, no repair} & 70.8\% & 41.7\% & \textbf{83.3\%} \\
\textbf{Staging} & 0.099 & 0.804 & \textbf{0.890} \\
\textbf{Scenes $\geq0.825$} & 0.0\% & 65.0\% & \textbf{80.0\%} \\
\bottomrule
\end{tabular}
\end{table}

\textbf{Reranking recovers everything sampling broke except action
fidelity, at the cost of $4\times$ inference} (a forward pass measured at
2.7~s): \texttt{valid} reaches 94.0\% in distribution and 83.3\% out of
it, both the highest this project has recorded, against v1's 68.7\%
baseline. Staging is essentially untouched ($0.875 \rightarrow 0.879$),
which is the point of tie-breaking on draw order rather than staging
score --- the rerank selects structurally clean scenes, and they happen
to stage just as well, rather than selecting for staging directly.
\textbf{What rerank does not recover is \texttt{actions\_per\_character}}
(85.7\% under beam4 against 72.7\% here) --- expected, since the rerank
criterion is structural validity, orthogonal to whether the right
character waved. \textbf{There is no single best configuration}:
beam4+colour-processor maximizes in-distribution obedience (ALL 57.1\%,
actions 85.7\%, staging 0.182); sampled+rerank+colour-processor maximizes
composed appearance (staging 0.879, valid 94.0\%, actions 72.7\%) --- the
two differ by 0.70 of the staging metric on identical weights, and the
paper states which is used, every time.

\subsection{The deliverables}
\label{sec:deliverables-a23}

\begin{verbatim}
python model/scene_infer.py \
  "two friends meet in the park, one waves, then they kick a ball around" \
  --ckpt model/scene_ckpt3 --max-in 160 --max-out 256 \
  --sample --unique-colours --rerank 4 --seed 23 -o scenes/a23_generated.scene
\end{verbatim}

The same prompt Attempt 22 used, directly comparable:

\begin{table}[h]
\centering
\caption{Attempt 22 vs.\ Attempt 23 on the same prompt.
(Source: \texttt{ATTEMPTS.md} Attempt 23 \secsym13.)}
\label{tab:a22vsa23}
\begin{tabular}{@{}L{4.3cm}L{4cm}L{4cm}@{}}
\toprule
 & A22 \texttt{a22\_generated.mp4} & A23 \texttt{a23\_generated.mp4} \\
\midrule
Cast positions & 9.6 and 9.2 & \textbf{7.9 and 3.1} \\
Minimum separation & 0.4 & \textbf{1.385} \\
Span of the usable stage & 15.6\% & \textbf{29.6\%} \\
Occlusion fraction & 0.344 & \textbf{0.000} \\
\textbf{Staging} & \textbf{0.332} & \textbf{0.886} \\
Critic warp error ($\leq$0.06) & 0.0022 & 0.0014 \\
Critic motion coverage ($\geq$0.2\%) & 2.03\% & 1.49\% \\
Identity worst & 1.00 & 0.9999 \\
CLIPSIM & 0.297 & 0.3177 \\
\textbf{Critic verdict} & USABLE & \textbf{USABLE} \\
Cast audit & 2/2 visible, 55/55 frames & \textbf{2/2 visible, 82/82 frames} \\
\bottomrule
\end{tabular}
\end{table}

\textbf{Staging 0.886 puts a model-written scene above
\texttt{street\_relay.scene} (0.825) for the first time in this
project.} Motion coverage is lower (1.49\% vs.\ 2.03\%) and this is not a
hidden regression: the clip is longer and the characters stand further
apart, so less of the frame changes per unit time, and it still clears
the critic's 0.2\% floor seven-fold.

\textbf{Second deliverable: a cast of six}, unreachable by any earlier
model in this project. From ``six friends spread out on a night street,
all of them dancing, then two of them kick a ball around'' $\rightarrow$
\texttt{scenes/a23\_six.scene} $\rightarrow$ \texttt{out/a23\_six.mp4}:
11.3~s, 339 frames, \textbf{6 characters}, 2 props, 16 timeline events,
zero structural problems, zero repairs.

\begin{table}[h]
\centering
\caption{Six-character deliverable. (Source: \texttt{ATTEMPTS.md} Attempt 23 \secsym13.)}
\label{tab:sixcast}
\begin{tabular}{@{}lr@{}}
\toprule
 & Value \\
\midrule
Staging & \textbf{0.903} \\
Span of the usable stage & \textbf{89.8\%} \\
Critic warp error & 0.0077 \\
\textbf{Critic motion coverage} & \textbf{6.82\%} \\
Identity worst & 0.9993 \\
\textbf{Critic verdict} & \textbf{USABLE} \\
Cast audit & \textbf{6/6 visible in 68/68 sampled frames} \\
\bottomrule
\end{tabular}
\end{table}

\textbf{6.82\% motion coverage is the highest of any clip in this
project} (\texttt{out/a21\_park.mp4}, the hand-written reference, is
5.32\%), from six independently animated characters composed in one shot
from one English sentence. Both critic runs used
\texttt{tools.gpuguard.Guard} at 50\% duty (0 cooldowns); generation,
rendering and every model in this attempt are CPU-only --- the card is
touched only by the critic's YOLO pass.

\subsection{Scoring every prediction, including the ones that were wrong}
\label{sec:scoringpredictions}

\begin{table}[h]
\centering
\footnotesize
\caption{Every RESEARCH2 \secsym7 prediction, scored honestly.
(Source: \texttt{ATTEMPTS.md} Attempt 23 \secsym14.)}
\label{tab:scoringpredictions}
\begin{tabular}{@{}cL{5cm}L{7cm}@{}}
\toprule
\# & Prediction (RESEARCH2 \secsym7) & Verdict \\
\midrule
1 & Colour processor: duplicates $71/150\rightarrow0$, nothing else
  changes & \textbf{half right} --- $0/150$ and $0/24$ exactly, but
  masking shifts the beam path (v2 \texttt{ALL} $86.0\rightarrow88.0$,
  v3 \texttt{parse} $98.7\rightarrow98.0$) \\
2 & Spacing pass: minimum pairwise distance rises to an enforced floor &
  \textbf{falsified as stated} --- true of declared positions ($0.157
  \rightarrow 1.217$), false of the clip ($0.069 \rightarrow 0.348$),
  because the pass never touches walk destinations \\
3a & Staged position sampling: v3 stages well above v2 without any pass
  & \textbf{right, invisible under beam search} --- 55 distinct $x$
  under sampling vs.\ v2's 20; \emph{worse} under beam4 (0.182 vs.\
  0.363) because a wider distribution flattens the decoding peak
  further \\
3b & Spatial language: a measurable rate of obeying a stated instruction
  & \textbf{falsified} --- $+2.7$ points over a mismatched-pair control
  in distribution, $-2.0$ under sampling; the 44.3\% headline is the
  base rate \\
4 & Cast range 1--8: 6--8 read correctly, 9 saturates & \textbf{confirmed
  in both halves} --- 15/15 inside range, cast of 6 every time at 9 \\
5 & \texttt{d\_model} 512 rules out an embedding-width bottleneck &
  \textbf{confirmed} against the actual \texttt{config.json} \\
--- & \emph{(not predicted, found)} & \textbf{the position collapse is
  mostly a decoding artefact} --- $0.188\rightarrow0.862$ on identical
  weights \\
\bottomrule
\end{tabular}
\end{table}

Two of five stated predictions were falsified, in whole or in part, and
are reported that way rather than reframed as successes after the fact.

\subsection{Where this attempt and RESEARCH2 disagree}
\label{sec:disagree}

RESEARCH2 ranks the deterministic spacing pass second among its
recommendations and does not consider the decoding strategy as a
candidate cause at all; its \secsym2 diagnosis of ``conditional
mode-collapse'' was the right read of the \emph{distribution} and, absent
a decoding-strategy control, the wrong read of the \emph{symptom} --- the
measured cause is that beam search reports the mode of a distribution
that turned out to be adequate, and the fix was one argument to
\texttt{generate()}, not a factor graph. We state this as a limitation of
the diagnostic tools available when the survey was written, not a
criticism of the survey: nothing in the project's log at that point had
measured staging at all, so nothing distinguished ``the model collapsed''
from ``the search collapsed it.'' Separately, RESEARCH2 \secsym4's
recommendation to hand-write the one constraint needed and hold
Outlines/Guidance as a fallback was followed and, per \secsym11.9, the
fallback is now worth reconsidering --- once sampling replaces beam
search, unconstrained decoding emits tokens outside the DSL grammar
entirely, which a single hand-rolled processor does not cover.

\subsection{What is honestly not solved}
\label{sec:notsolved-a23}

The model does not obey spatial instructions above the base rate
(\secsym11.11). Cast size saturates at 6 when asked for 9 --- a
relocated wall, not a removed one, consistent with a known general
seq2seq limitation (\secsym11.10). Beam search and sampling are 0.70
apart on staging and 13 points apart on action fidelity on identical
weights, and no single decoding configuration is best across every
metric (\secsym11.9, \secsym11.12) --- the pipeline currently makes the
user choose. \texttt{valid} under sampling is carried by rerank (4
forward passes to launder one), not by the model directly. The spacing
pass does not touch the timeline, so a scripted \texttt{walk to} can
still undo it. The staging metric itself is calibrated on $n=13$ clips
(2 hand-written, 11 model-written) and is measurably weaker above 4
characters, where a single bad pair is diluted across many pairwise
terms --- \texttt{a23\_six} has three characters within 0.2 units of
each other and still scores an occlusion fraction of only 0.067.
Inference truncates large casts (v3 parse 98.7\% is entirely 7--8-
character scenes exceeding \texttt{--max-out 256}, a cap, not a
capability limit). This is one seed on one architecture; the decoder
result is far too large to be seed noise, but the 24- and 12-prompt
out-of-distribution sets make per-group numbers indicative, not tight.
Every limitation Attempts 21 and 22 listed that this attempt did not
touch still stands: foot slide, no collision, actions cannot blend
across a clip boundary, canonical character names, unmeasured pacing,
and CLIPSIM's inability to judge this domain at all.

\section{Answering the Project's Own Indictment}
\label{sec:indictment}

The project's standing self-criticism, stated identically at the end of
Era 1 and again at the end of Attempt 21, is: \textbf{``the models never
beat the script.''} Attempt 22 answers it directly, and both halves of
the answer are reported, not just the favorable one.

\begin{table}[h]
\centering
\caption{Hand-written vs.\ model-written, Attempt 21 vs.\ Attempt 22.
(Source: \texttt{ATTEMPTS.md} Attempt 22, \secsym9, ``Attempt 4's
indictment, answered with a number.'')}
\label{tab:indictment}
\begin{tabular}{@{}L{3.2cm}L{4.7cm}L{5.2cm}@{}}
\toprule
 & Hand-written (Attempt 21) & Model-written (Attempt 22) \\
\midrule
Clip & \texttt{out/a21\_park.mp4} & \texttt{out/a22\_generated.mp4} \\
Author & a human, by eye, over several passes & t5-small, 60M params, one
  forward pass \\
Verdict & \textbf{USABLE} & \textbf{USABLE} \\
Warp error & 0.0055 & \textbf{0.0022} \\
Motion coverage & \textbf{5.32\%} & 2.03\% \\
Cast & 3 & 2 \\
Props & 7 & 2 \\
Timeline events & 8 & 4 \\
Duration & 12.5~s & 9.1~s \\
Staging & positions chosen by hand, two moved after looking at the
  frames & two characters at $x{=}9.6$ and $x{=}9.2$, stacked at the
  right edge of a 10-unit frame \\
\bottomrule
\end{tabular}
\end{table}

\textbf{The straight answer: no. The model did not beat the hand-written
script --- it reached the same verdict with a thinner scene.} Both clips
are USABLE and the generated one has a \emph{lower} warp error, but warp
error rewards stillness, and the honest reading of 0.0022 against 0.0055
is ``less is happening,'' not ``better.'' The row that matters is motion
coverage (2.03\% against 5.32\%) and the row under it: the human's scene
has three characters, seven props and eight events against the model's
two, two and four, and the human's characters are spread across the
frame because a human looked at a frame and moved them, while the
model's are stacked at one edge because the prompt never said where
anyone should stand.

\textbf{What has changed, and it is not nothing:} the sentence ``100\%
hand-written'' is now false --- a natural-language prompt the model had
never seen produced a valid, unrepaired scene script that rendered to a
USABLE video with no human in the loop. 90.7\% of held-out prompts are
satisfied completely (v1, in-distribution; 86.0\% for v2), with 100\%
parse and 100\% render, and 0\% exact-match --- the model is not
reproducing scripts, it is writing new ones. 9 of 10 generated clips
clear the critic, and the tenth fails a motion floor for correctly doing
what its prompt asked.

\textbf{What has not changed:} the best-\emph{looking} output of this
project is still \texttt{out/a21\_park.mp4}, and it is still
hand-written. The model now writes scripts that are structurally correct
and semantically obedient; it does not yet write scripts that are well
\emph{staged}, and staging is what made the Attempt 21 clip look like a
scene rather than merely be a correct one. The project's indictment is no
longer true as originally written, and its spirit still stands.

\textbf{Attempt 23 updates this comparison on exactly one axis, and it is
worth stating precisely rather than folding it back into a cleaner
story.} With a corrected decoder (sampling plus rerank, not the beam
search this section's table used), the same prompt now produces
\texttt{out/a23\_generated.mp4} at staging \textbf{0.886} --- above
\texttt{street\_relay.scene} (0.825), the weaker of the project's two
hand-written reference clips, for the first time in this project's
history (\secsym11.13). This is a real change to ``the model does not
stage well,'' which Attempt 22 treated as settled. It is not a change to
the rest of the table above: the Attempt 23 deliverable still has 2
characters, 2 props and 4 events against the hand-written park scene's
3, 7 and 8, and --- the sharper point --- the same attempt's own control
shows the model's apparent obedience to stated spatial instructions is
within noise of a mismatched-prompt baseline ($+2.7$ points,
\secsym11.11). \textbf{The model now composes a well-spaced scene; it
does not yet take direction about where to place anyone.} Attempt 23's
own closing sentence states this distinction as precisely as we can:
``this model composes, and does not yet take direction.''

\section{Limitations}
\label{sec:limitations}

\textbf{Arc one (diffusion route) is closed, not merely unfinished.} No
configuration of SD1.5 + ControlNet + AnimateDiff tried in this project
reaches the evaluation harness's USABLE verdict; the mechanism that helps
flicker measurably costs pose obedience, and the mechanism that helps
identity measurably costs a little flicker back. We do not expect further
tuning of this exact stack to resolve this on the reported hardware ---
the finding is that the objectives trade against each other, not merely
that none was tuned enough.

\textbf{Arc two (symbolic composition) does not yet compose style or
take spatial direction.} \texttt{scenescript.py} produces line-art stick
figures, not styled anime characters --- this project has not combined
arc two's composition machinery with arc one's diffusion styling, and
doing so is unattempted. Cast size saturates at 9 for the current model
and training range (1--8) --- a wall relocated by exactly the width of
the range extension, not removed, and consistent with a documented
general limit of sequence-to-sequence counting/extrapolation (RESEARCH2
\secsym3) rather than a defect specific to this system (\secsym11.10).
Duplicate character colours are fully eliminated by a
constrained-decoding logits processor ($0/150$, $0/24$, 0 across 348
generated scenes, \secsym11.8) --- this one wall is closed, not merely
patched. \textbf{The model does not obey stated spatial instructions
above the rate a mismatched control achieves} ($+2.7$ points in
distribution, \secsym11.11) --- every staging improvement in this paper
comes from the decoding strategy and the position sampler, not from the
model reading ``leftmost'' or ``close together.''

\textbf{Getting a well-staged scene costs something, and the cost is
stated plainly.} Sampling instead of beam search buys 0.693 of the
staging metric at a cost of 13--20 points of in-distribution obedience
(\texttt{actions\_per\_character} $85.7\% \rightarrow 72.5\%$,
\texttt{layout} $45.8\% \rightarrow 25.3\%$), concentrated exactly in the
fields the prompt specifies; sample-and-rerank recovers structural
validity ($60.7\% \rightarrow 94.0\%$) but not action fidelity (72.7\%,
still 13 points below beam search's 85.7\%, \secsym11.9, \secsym11.12).
There is no single decoding configuration that is best on every axis, and
the pipeline currently requires choosing one per use case rather than
offering one that dominates.

\textbf{The staging metric itself is a project-defined instrument,
calibrated on a small sample.} Its separation between hand-written and
model-written scenes is complete ($n=13$: 2 hand-written, 11
model-written, no overlap), but $n=13$ is small, two of its five
candidate sub-terms were discarded during calibration for disagreeing
with the reference clips (\secsym11.2), and it is measurably weaker
above four characters, where a single close pair is diluted across many
pairwise terms (\secsym11.16). It has not been validated against human
judgment of what ``well staged'' means, which is a real gap (see the
user-study limitation below).

\textbf{The deconfounding fix (Attempt 22) trades 4 points of
in-distribution accuracy for 19.5 points of out-of-distribution
improvement, and this trade is reported, not minimized.} The remaining
23.5-point gap after deconfounding is not fully understood --- the
\texttt{valid} rate's own regression ($68.7\% \rightarrow 40.0\%$) shows
the fix is not free even on the axis it targeted.

\textbf{This paper reports no user study, and no comparison against a
competing system.} All quality judgments in this paper are made by this
project's own automated critic (\secsym4) and staging metric
(\secsym11.2), both built and calibrated internally, and by close reading
of individual clips and generated strings --- never by an independent
human rater blind to which clip came from which condition, and never
against another published script-to-animation system run on the same
prompts. RESEARCH2 \secsym5 found no other system in this specific niche
(sub-100M parameters, CPU-trainable, script-to-multi-character-animation)
to compare against as far as that survey's search found, which explains
the absence of a system-level baseline but does not substitute for one; a
human preference study comparing hand-written, beam-search, sampled, and
reranked outputs is the most direct way to close this gap and was not run
here.

\textbf{Every rate in Attempt 23 beyond the headline decode probe is
measured on a fixed set of 150 in-distribution and 24
out-of-distribution prompts, on one seed, one architecture (t5-small),
one training run per checkpoint.} The decoder result itself (0.188 vs.\
0.862 staging on identical weights) is far too large a gap to be seed
noise, but the 24-prompt (and, for the spatial-language control,
12-prompt) out-of-distribution sets make per-group numbers indicative
rather than tight, and no result in this paper has been replicated
across multiple training seeds.

\textbf{Future work, stated with the reasoning rather than as an
unexplained gap:} a DETR-style Hungarian set-decoder head
\cite{carion2020detr} for the cast section is, per RESEARCH2 \secsym4,
the one surveyed technique that structurally addresses staging and
duplicate colours together --- $N$ cast slots predicted with mutual
attention would make a repeated colour visible to the model during
generation itself, and a repulsion term over the predicted positions
would attack staging in the same head. It is the only technique surveyed
requiring genuine architecture work (a new decoder head, a new
order-invariant matching loss) against this project's ``minutes of CPU,
no architecture change'' retraining budget, and it was deliberately not
built: a one-argument decoding change (sampling) and a cheap rerank pass
closed most of the measured gap without it (\secsym11.6, \secsym11.12),
so it is no longer the obvious next move, but it remains the principled
fix if extending the constrained-decoding approach to the full DSL
grammar (also future work, per \secsym11.9) does not close what remains
of the \texttt{valid} and action-fidelity gaps under sampling.

\textbf{Pose conditioning obedience remains partial after Attempt 19's
corrections and then gets strictly worse under the Attempt 20 motion
module} (\secsym7.1, \secsym8.4) --- an obeyed guide pose falls from 1 of
4 inspected frames to 0 of 8 once temporal attention is added.

\textbf{Identity preservation via IP-Adapter (\secsym8.3) is measurably
improved but only on a narrow, stated axis} --- colour scheme, not
structural identity, per the critic's own HSV-histogram embedding and the
anchor run's contact sheet.

\textbf{No off-the-shelf pose or person detector can currently score
this project's styled, symbolic, or multi-character output} (\secsym5.3,
\secsym7.3, \secsym9.5) --- a real, unresolved measurement gap in the
available tooling for 2D line art of any kind, not specific to any one
clip.

\textbf{Content-correctness evaluation remains asymmetric across the
project's output paths} (\secsym4.5) --- exact for the symbolic AniSVG
path, absent by construction for the diffusion-rendered path, and not
yet extended to Attempt 21/22's composed scene scripts even though they
are, like AniSVG, a declarative source that could in principle be
checked exactly.

\textbf{FVD, content-debiased FVD, JEDi, and FVMD are not computed
anywhere in this project} because the clip counts available (dozens per
condition) do not meet the sample sizes the published literature uses
for a stable Fr\'echet estimate; this paper reports warping-error-style
and flicker-ratio metrics instead, per the survey's own recommendation.

\textbf{Several numbers cited in Related Work (\secsym2) are marked
unverified in the source survey itself} and are reported only as
context, not as claims this project has confirmed locally.

\textbf{Reproducibility caveat:} all GPU experiments ran under a 50\%
duty-cycle guard for thermal reasons; numbers measured on undamaged or
unthrottled hardware are not expected to be directly comparable on
wall-clock terms, though the dimensionless ratios (flicker, PCK,
prompt-match rates, staging) should reproduce. Attempts 21--23 use no GPU
except the critic's own YOLO pass and are therefore not subject to this
caveat for their own generation and training --- a property we consider
part of their result, not a footnote (\secsym9, \secsym10.2,
\secsym11.5).

\section{Conclusion}
\label{sec:conclusion}

Twenty-three logged attempts split this project's account into two arcs,
and both are reported honestly rather than as one continuous march toward
a working system. The first arc closes: Stable Diffusion~1.5, ControlNet
and AnimateDiff, on this hardware and with the corrections this project
applied, reduce flicker for the first time in the project's history
($36.2\times \rightarrow 25.0\times$) and still cannot clear the
evaluation harness's own bar, and the mechanism that helps flicker
measurably costs pose obedience while the mechanism that helps identity
measurably costs flicker back. That interference --- \textbf{flicker,
identity and pose fidelity behave as separable defects on this hardware
and model family, not sub-problems of one shared deficiency} --- is this
paper's first central finding, and it rests on evidence spanning three
attempts: Attempt 19 showed pose obedience and flicker orthogonal even
with no temporal mechanism present; Attempt 20 supplied the temporal
mechanism and traded it directly for pose; and underlying both, Attempt
19's proportion finding shows the very choice that makes a character read
as anime is the choice that removes its pose guide from a human-pose
model's training distribution.

The second arc opens, and gets further than the first, in three steps. A
CPU-only symbolic compositor (Attempt 21) produces this project's first
multi-character animation and clears the evaluation harness's temporal
gates by a wide margin. A 60M-parameter model (Attempt 22) trained in
under nine minutes on a CPU then learns to write scripts for it ---
generalizing at the format level essentially perfectly and failing
compositionally in one diagnosable way. That diagnosis --- tracing a
65.7-point generalization gap to a spurious correlation the project put
into its own synthesized training data, confirming it with a controlled
probe, fixing it at the data level, and reporting the resulting trade
honestly --- is this paper's second central finding, and we hold it up as
a general methodological point: \textbf{a small model trained on
synthesized data can look incapable of a whole class of reasoning when
the truer fault is what its training data ever asked it to learn}, and
this is distinguishable, by a controlled probe, from a genuine capability
limit.

\textbf{A third attempt (Attempt 23) then does the same thing to the
project's own diagnosis of itself, and this is the paper's third central
finding.} Attempt 22 named staging as the model's biggest remaining
defect; decoding the identical checkpoint on identical prompts under
sampling instead of beam search moves staging from 0.188 to 0.862 on one
checkpoint and 0.333 to 0.845 on another, and the number of distinct
positions emitted across thirty scenes goes from 3 to 55. The defect the
project had attributed to the model was mostly a property of the search.
A companion ablation shows precisely what would have been believed
without this check: a deterministic spacing pass looks, on its own, like
a successful fix (staging $0.178 \rightarrow 0.587$), and only comparison
against decoding alone ($0.178 \rightarrow 0.862$, no pass) reveals that
beam search was suppressing an already-adequate distribution the whole
time --- \emph{good engineering, bad science claim}, stated as the
project states it. This completes a pattern rather than standing alone:
this project has now overturned its own first explanation of an apparent
model failure three times --- a topology bug that was fixed and changed
nothing until a proportion mismatch was found by controlled ablation
(Attempt 19), a ``cannot count'' verdict that a clause-count probe traced
to a 99.9\% data confound (Attempt 22), and now a ``cannot stage''
verdict traced to a decoding default (Attempt 23) --- and we report the
pattern of self-correction itself as a contribution, alongside the
individual findings it produced.

Running through all three arcs is a fourth contribution we take as
seriously as any of them: this project's own evaluation harness and
decoding pipeline produced four false or unmeasurable verdicts along the
way --- a content-correctness metric that could not separate correct
geometry from known-bad geometry, an identity metric that silently mixed
two measurement bases and cost one real clip its verdict twice, person
detectors blind to every stick figure this project has drawn, and a
decoder that reported the mode of a fine distribution as if it were the
distribution itself --- each caught and corrected in the open. The
underlying lesson generalizes beyond any one of them: a metric's or a
decoding default's provenance bounds what it can be trusted to measure or
produce, and this bound must be checked against a control, never assumed.

We close by answering the project's own standing indictment precisely,
not by resolving it in either direction. The model still does not beat
the hand-written script on richness: Attempt 22's generated scene is
thinner than the human-authored one (2 characters, 2 props, 4 events
against 3, 7, 8), and that comparison is unchanged by anything in this
paper. What has changed, twice, is real and stated exactly. First: the
claim that this project's best output is ``100\% hand-written'' is false,
because an unseen sentence produces a valid, unrepaired script that
renders to a video the evaluation harness calls USABLE, with no human in
the loop (Attempt 22). Second: with a corrected decoder, a model-written
scene now \emph{stages} better than the weaker of the project's two
hand-written reference clips ($0.886$ vs.\ $0.825$) --- composition,
specifically --- while the same attempt's own control shows the model
still does not reliably obey a stated spatial instruction above the rate
a mismatched control achieves (Attempt 23). The model composes; it does
not yet take direction. We take the position stated in the introduction:
on hardware this constrained, a correctly diagnosed and measured result
--- a closed route, an open one, a defect traced to its actual cause, or
a defect traced away from where the project first placed it --- is a
more useful contribution than an unmeasured or overclaimed one, and we
have tried throughout, including about our own project's prior
conclusions, to write this paper that way.

% ATD-12K and MDM (Tevet et al.) are listed in paper/PAPER.md's flat
% References section (carried over from paper/RESEARCH.md's dataset and
% motion-model survey) without an inline in-text citation in the
% condensed Related Work prose; \nocite them so they still appear in the
% bibliography, matching PAPER.md's own reference list exactly.
\nocite{atd12k,tevet2022mdm}

\bibliographystyle{plain}
\bibliography{refs}

\end{document}
